\documentclass[11pt]{article}
\usepackage[T1]{fontenc}
\usepackage{lmodern,microtype,amsmath,amssymb,amsthm,mathtools,mathrsfs,tikz-cd}
\usepackage[margin=1in]{geometry}
\usepackage{enumitem,xurl,hyperref,needspace}
\hypersetup{colorlinks=true,linkcolor=blue,citecolor=blue,urlcolor=blue,
pdftitle={Definability and incompleteness in recursion categories},pdfsubject={Version 16, 27 October 2026}}
\newtheorem{theorem}{Theorem}[section]
\newtheorem{proposition}[theorem]{Proposition}
\newtheorem{lemma}[theorem]{Lemma}
\newtheorem{corollary}[theorem]{Corollary}
\theoremstyle{definition}
\newtheorem{definition}[theorem]{Definition}
\theoremstyle{remark}
\newtheorem{remark}[theorem]{Remark}
\title{Definability and incompleteness in recursion categories}
\author{}
\date{27 October 2026\\Version 16}
\begin{document}
\maketitle
\begin{abstract}
We prove a categorical essential incompleteness theorem for recursion
categories equipped with additional structure expressing definability.
The proof establishes a categorical analogue of recursive inseparability.
In the numerical interpretation, the definability hypothesis is equivalent
to strong numeralwise definability of every recursive set in a consistent
arithmetical theory $T$ with effective syntax and numeral substitution.
The conclusion is essential incompleteness: every consistent c.e.\ extension
of $T$ is incomplete. This gives a categorical version of one of the
generalized incompleteness statements proposed by Di Paola and Heller.

Let $\mathscr V$ be a finitely complete category. A definability algebra
consists of a $\mathscr V$-enriched category $\mathcal K$ whose underlying
ordinary category $\mathcal K_0$ is a recursion category, together with
predicate objects, compatible restriction and inverse-image operations,
embeddings of restriction idempotents, and internal indexing and
representation operations in $\mathscr V$. The representation operation
assigns a total representing arrow to every complemented restriction
idempotent on a fixed object, relative to a fixed total arrow and two
fixed disjoint predicates.

If $\mathcal K_0$ is nontrivial, two disjoint predicates that arise
from restriction idempotents and contain the two fixed predicates,
respectively, cannot have join equal to the greatest element of their
predicate lattice. We formulate categorical essential incompleteness
using quotients of internal Boolean algebras and derive it from this
result.
\end{abstract}

\section{Introduction}
\label{sec:introduction}

Di Paola and Heller sought an algebraic formulation of generalized
incompleteness in which the recursive constructions and diagonal
argument are developed within recursion categories
\cite[pp.~594--596]{DPH}.
One of their formulations asserts that definability of every
recursive set in a consistent arithmetical theory implies recursive
inseparability of its theorem and refutation sets. Essential
incompleteness follows: every consistent c.e. extension is incomplete.
We prove a categorical form of this implication and a categorical
essential incompleteness theorem. The additional structure expressing
definability is specified by a definability algebra.

We write $\mathcal K$ for an enriched category and $\mathcal K_0$
for its underlying ordinary category. We use Heller's convention:
a recursion category $\mathcal K_0$ is a prodominical
category in which all objects are isomorphic and which admits a
Turing morphism \cite[Sections~5 and~9]{Heller}.
For an arrow $f:X\to Y$, write $\bar f:X\to X$ for its
restriction. An endomorphism $e:X\to X$ satisfying $\bar e=e$
is a \emph{restriction idempotent}
\cite[Section~2.1]{RestrictionI}.
These are the morphisms called \emph{domains} in Di Paola and
Heller's terminology \cite[Definition~3.1 and Lemma~3.2]{DPH}.
We use \emph{restriction idempotent} throughout.
A total arrow satisfies $\bar f=1_X$;
the total arrows form $\operatorname{Tot}(\mathcal K_0)$.
We write $\otimes$ for binary restriction products and use the
composition convention $gf=g\circ f$: first $f$, then $g$.

Fix a finitely complete category $\mathscr V$ and a
$\mathscr V$-enriched category $\mathcal K$ whose underlying
ordinary category $\mathcal K_0$ is a recursion category.
A definability algebra equips $\mathcal K$ with restriction
operations, predicate objects, inverse-image operations, and
embeddings of the objects of restriction idempotents into the
predicate objects. It also includes internal indexing and
representation operations in $\mathscr V$. The representation
operation assigns a total representing arrow to every complemented
restriction idempotent on a fixed object, relative to a fixed total
arrow and two fixed disjoint predicates. Its source is a pullback
in $\mathscr V$ expressing complementarity. A second pullback in
$\mathscr V$ expresses the representation equations, and a section
of its projection supplies the operation.

Theorem~\ref{thm:internal-inseparability} proves that, in a nontrivial
recursion category, no complemented restriction idempotent separates
the two distinguished predicates. Suppose the predicate lattices
are realized as subobject lattices in a regular extensive category,
and the distinguished predicates are the fibers over the bounds
of a Boolean presentation. Every consistent presentation induced
by a Boolean quotient is incomplete if its two fibers are images
of restriction idempotents
(Theorem~\ref{thm:categorical-essential-incompleteness}).
In the numerical interpretation, the representation condition is
equivalent to strong numeralwise definability of every recursive
set, and the categorical inseparability conclusion becomes
recursive inseparability of the theorem and refutation sets
(Theorem~\ref{thm:integration-essential-equivalence}).
The Boolean quotient conclusion is equivalent to essential
incompleteness (Proposition~\ref{prop:numerical-boolean-extensions}).
Consistent extensions of Peano arithmetic provide the numerical
examples.

\subsection{Related literature}
\label{sec:related-literature}

Lawvere's 1969 categorical diagonal and incompleteness theorems
use representation by global elements
\cite[Theorems~2.2 and~3.3]{Lawvere}.
Di Paola and Heller's 1987 program instead asks for the recursive
constructions and generalized incompleteness statements within
recursion categories. The numerical interpretation of our
inseparability theorem recovers their formulation~(3)
\cite[pp.~594--596]{DPH}.
They reject diagonalizable algebras for this purpose: first
incompleteness reduces to the assertion that a nontrivial algebra
has more than two elements, while the axiom formalizing L\"ob's
theorem largely incorporates second incompleteness.
Their proposed development uses intrinsic recursion theory to
obtain the generalized incompleteness statements.

Di Paola and Heller also construct recursion categories of free
actions of finite groups
\cite[Section~7.9(3), Theorem~7.9.9]{DPH}.
The equivariant recursion category used here is the one-object
form of that construction over the numerical recursion category.
Heller's 1990 existence theorem supplies recursion categories
beyond the numerical prototype
\cite[Theorem~9.2 and Section~10]{Heller}.
His formulation adopts Rosolini's P-categories and uses the
prodominical condition \cite[Introduction and Section~5]{Heller}.

Cockett and Hofstra's 2008 Turing categories replace the isotype
condition by a Turing object
\cite[Sections~3.1 and~3.6]{TuringCategories}.
Cockett's March 2010 lectures attribute a recursive-inseparability
theorem in this setting to Lengyel
\cite[slides~33--34]{CockettSlides}.
His 2014 notes state the result for discrete Turing categories
with finite joins and give an incompleteness theorem with a
provability predicate
\cite[Section~7.1 and Theorems~7.8 and~7.14]{CockettNotes}.
The categorical hypothesis used here is representation of
complemented restriction idempotents by fixed disjoint predicates;
Theorem~\ref{thm:integration-essential-equivalence} identifies
its numerical meaning.

Ramesh's 2023 introspective theories provide categorical fixed-point
theorems \cite[Chapters~2--4]{Ramesh}.
We construct a relative free introspective extension that preserves
a definability algebra, including its representation operation.
For the arithmetic introspective structures, quotation and
consistency supply the numerical representation operation after
standard interpretation. This construction uses actual totality
of a decision procedure, without requiring its totality to be
provable in the arithmetical theory.

The numerical definability algebra is enriched over $\mathbf{Set}$;
its Boolean presentations belong to the c.e. and full categories
of partial equivalence relations. On each of these categories,
every introspective functor is naturally isomorphic to a constant
functor with terminal value. Conversely, the PA syntactic base
admits nonconstant introspection, but its enrichment of programs
modulo uniformly PA-provable graph equality admits no internal
representation operation. Among the base and enrichment pairs
specified in Appendix~\ref{sec:introspection-representation-appendix},
neither nonconstant introspection
on the base nor extension to a definability algebra implies the other.

Definability algebras were introduced by adjoining chosen total
representing arrows to the recursion-category structure
\cite[Sections~3--4]{DefinabilityAlgebras}.
The present construction internalizes the representing operation
and allows general predicate objects. Taking the objects of
restriction idempotents as predicate objects recovers the original
algebras. Their inclusion into the general predicate structures
is fully faithful; the numerical interpretation shows that it
need not be essentially surjective.

\subsection{Organization}

Section~\ref{sec:integration} defines the categorical structure,
proves inseparability, and establishes categorical essential
incompleteness for Boolean presentations.
Section~\ref{sec:integration-numerical} proves the numerical
realization and its exact-completion comparison.
Section~\ref{sec:equivariant-model} gives the equivariant example,
including its Turing morphism and Boolean presentation.
Section~\ref{sec:completion-introspection} constructs extensions
with restriction terminal objects, exact completions, the relative
free introspective extension,
and the comparisons with nonconstant introspection. It also
derives the numerical definability structure from arithmetic
introspection.
Appendix~\ref{sec:introspection-representation-appendix} proves
the obstruction to an internal representation operation over
the PA syntactic base.

\section{Definability algebras and categorical incompleteness}
\label{sec:integration}

\subsection{Enriched categories and underlying recursion categories}
\label{sec:definability-algebras}

Fix a finitely complete category $\mathscr V$ with terminal object $1$.
All limits and structure arrows in this subsection belong to $\mathscr V$.
We use categories enriched over the cartesian monoidal category
$(\mathscr V,\times,1)$ in the sense of
\cite[Sections~1.2--1.3]{Kelly}.
Let $\mathcal K$ be such an enriched category, with hom-objects
$H(X,Y)$ in $\mathscr V$, identity
arrows $1\to H(X,X)$, and composition arrows
\[
 H(Y,Z)\times H(X,Y)\longrightarrow H(X,Z)
\]
satisfying the category equations. Restriction is given by arrows
$r_{X,Y}:H(X,Y)\to H(X,X)$ satisfying the restriction equations.
Assume chosen binary restriction products and restriction zeros, with their operations given by arrows of $\mathscr V$.
The restriction and product equations are those of
\cite[Section~2.1]{RestrictionI} and
\cite[Section~4.1]{RestrictionIII}, interpreted as equations
of morphisms in $\mathscr V$.
Write $\Gamma=\mathscr V(1,-):\mathscr V\to\mathbf{Set}$.
The underlying ordinary category $\mathcal K_0$ has the same
objects as $\mathcal K$ and hom-sets
\[
 \mathcal K_0(X,Y)=\Gamma H(X,Y)=\mathscr V(1,H(X,Y)).
\]
Applying $\Gamma$ to the enriched identities, composition,
restriction, and product operations gives the corresponding
operations in $\mathcal K_0$. Require $\mathcal K_0$ to be a
recursion category. For objects $X,Y$ of $\mathcal K$, an arrow
$f:X\to Y$ of $\mathcal K_0$ is represented by a point
$\widehat f:1\to H(X,Y)$ in $\mathscr V$. We distinguish these
notations when both types occur in the same construction.

The objects of restriction idempotents and total arrows are defined by
\[
\begin{tikzcd}
 E_X \arrow[r, hook, "i_X"] & H(X,X)
 \arrow[r, shift left, "r_{X,X}"]
 \arrow[r, shift right, swap, "1_{H(X,X)}"] & H(X,X),
\end{tikzcd}
\qquad
\begin{tikzcd}
 \mathsf T(X,Y) \arrow[r, hook] \arrow[d] & H(X,Y) \arrow[d, "r_{X,Y}"]\\
 1 \arrow[r, "1_X"'] & H(X,X),
\end{tikzcd}
\]
where the first diagram is an equalizer and the second a pullback.
For every object $A$ of $\mathscr V$, the equalizer identifies
morphisms $A\to E_X$ with morphisms $u:A\to H(X,X)$ satisfying
$r_{X,X}u=u$. In particular, preservation of finite limits by
$\Gamma$ gives
\[
 \Gamma E_X\cong E(X),\qquad
 \Gamma \mathsf T(X,Y)\cong\operatorname{Tot}(\mathcal K_0)(X,Y),
\]
where $E(X)$ consists of the restriction idempotents of $X$ in
$\mathcal K_0$. This identifies every restriction idempotent
$e:X\to X$ with a unique arrow $1\to E_X$.
The induced map $\Gamma i_X$ is the inclusion into the
endomorphism monoid $\mathcal K_0(X,X)$. The restriction equation
$e\bar e=e$ implies $e^2=e$ when $\bar e=e$.
Composition induces meet $\wedge:E_X\times E_X\to E_X$.
The restriction zero and identity give the lower and upper
bounds $\bot_{E_X},\top_{E_X}:1\to E_X$, respectively.
Equip $E_X$ with joins so that it is a bounded distributive lattice
object. Composition and restriction induce an inverse-image morphism
\[
 a^E_{X,Y}:H(X,Y)\times E_Y\longrightarrow E_X,
 \qquad (f,e)\longmapsto\overline{ef}.
\]
Require preservation of finite joins as equations of morphisms on
$H(X,Y)\times E_Y\times E_Y$. On $\mathsf T(X,Y)\times E_Y$ this
action also preserves the bounds. Expressions in variables here
and below abbreviate morphisms from the indicated finite products.

The chosen Turing morphism $U\otimes U\to U$ of $\mathcal K_0$
is represented by a global point $\tau:1\to H(U\otimes U,U)$
in $\mathscr V$. Composition with this point gives
\[
 \operatorname{ev}_\tau:\mathsf T(U,U)\longrightarrow H(U\otimes U,U),
 \qquad g\longmapsto\tau(g\otimes1_U).
\]
A chosen indexing arrow
$\Lambda:H(U\otimes U,U)\to \mathsf T(U,U)$ satisfies
$\operatorname{ev}_\tau\Lambda=1_{H(U\otimes U,U)}$.

\subsection{Predicate objects and the representation pullback}
\label{sec:generalized-definability-algebras}

\begin{definition}\label{def:generalized-definability-algebra}
A \emph{definability algebra} over $\mathscr V$ consists of the
enriched category $\mathcal K$ with the restriction, product,
lattice, and indexing structure above, together with the following
predicate and representation structure.
For each $X$, let $P_X$ be a bounded distributive lattice object.
Inverse image is given by arrows
\[
 a_{X,Y}:\mathsf T(X,Y)\times P_Y\longrightarrow P_X
\]
that preserve the lattice operations in the predicate argument and
satisfy the identity and contravariant composition equations
\[
 1_X^*a=a,\qquad (gf)^*a=f^*(g^*a).
\]
These are equations on $P_X$ and
$\mathsf T(Y,Z)\times \mathsf T(X,Y)\times P_Z$, respectively.
Let $j_X:E_X\rightarrowtail P_X$ be monomorphisms of bounded
distributive lattice objects, compatible with these inverse images.
Thus inverse image on restriction idempotents and on predicates satisfies
$a_{X,Y}(f,j_Y(e))=j_X(\overline{ef})$ as arrows on
$\mathsf T(X,Y)\times E_Y$.

An arrow $e:X\to X$ of $\mathcal K_0$ is represented
by a morphism $\widehat e:1\to H(X,X)$ in $\mathscr V$.
If $e$ is a restriction idempotent, the equation
$r_{X,X}\widehat e=\widehat e$ gives a unique factorization
$\widehat e=i_X\widetilde e$ through $E_X$.
Its associated predicate is $p_e=j_X\widetilde e$:
\[
\begin{tikzcd}[column sep=large,row sep=large]
 1 \arrow[r,"\widetilde e"] \arrow[dr,"\widehat e"']
   \arrow[rr,bend left=22,"p_e"]
 & E_X \arrow[r,"j_X"] \arrow[d,hook,"i_X"] & P_X \\
 & H(X,X). &
\end{tikzcd}
\]
Here $1$ is the terminal object of $\mathscr V$.
Thus predicates arising from restriction idempotents are exactly those arrows
$1\to P_X$ that factor through $j_X$.

Fix objects $U,S$ and a global point
$s:1\to \mathsf T(U\otimes U,S)$, representing a total arrow
$U\otimes U\to S$ of $\mathcal K_0$. Fix also predicates
$p_+,p_-:1\to P_S$ satisfying $p_+\wedge p_-=0$.
The object $C_U$ of pairs of complementary restriction idempotents is the pullback
\begin{equation}\label{eq:v8-complement-pullback}
\begin{tikzcd}
 C_U \arrow[r, hook, "d"] \arrow[d] & E_U\times E_U
 \arrow[d, "{\langle\wedge,\vee\rangle}"]\\
 1 \arrow[r, "{\langle\bot_{E_U},\top_{E_U}\rangle}"'] & E_U\times E_U.
\end{tikzcd}
\end{equation}
Write $d_+,d_-:C_U\to E_U$ for the two components of $d$.

Composition and tensor give an arrow
\[
 \chi:\mathsf T(U,U)\longrightarrow \mathsf T(U\otimes U,S),
 \qquad g\longmapsto s(g\otimes1_U).
\]
Put $W=P_{U\otimes U}$ and $G=\mathsf T(U,U)$. Define arrows
$\ell:G\to W\times W$ and $b:C_U\to W\times W$ by
\begin{align*}
 \ell_\pm&=a_{U\otimes U,S}\langle\chi,p_\pm!_G\rangle,\\
 b_\pm&=\pi_2^*j_Ud_\pm.
\end{align*}
Here $\pi_2^*:P_U\to W$ is inverse image under the second
projection $U\otimes U\to U$. Form the pullback
\begin{equation}\label{eq:v8-representation-pullback}
\begin{tikzcd}
 R \arrow[r, "r"] \arrow[d, "q_R"'] & G \arrow[d, "\ell"]\\
 C_U \arrow[r, "b"'] & W\times W.
\end{tikzcd}
\end{equation}
The representation structure is a section $\sigma:C_U\to R$ of $q_R$.
Put $\rho=r\sigma:C_U\to \mathsf T(U,U)$. The corresponding commutative
triangle is
\begin{equation}\label{eq:v8-representation-section}
\begin{tikzcd}
 C_U \arrow[rr, "\rho"] \arrow[dr, "b"'] && \mathsf T(U,U)
     \arrow[dl, "\ell"]\\
 & W\times W. &
\end{tikzcd}
\qquad q_R\sigma=1_{C_U},\qquad \rho=r\sigma.
\end{equation}
Equivalently, $\rho:C_U\to \mathsf T(U,U)$ satisfies
\begin{equation}\label{eq:generalized-algebra-representation}
 a_{U\otimes U,S}\langle\chi\rho,p_\pm!_{C_U}\rangle
   =\pi_2^*j_Ud_\pm.
\end{equation}
The pullback universal property makes these equations equivalent
to the existence of the section in
\eqref{eq:v8-representation-section}.
\end{definition}

\subsection{Morphisms and predicates arising from restriction idempotents}
\label{sec:definability-morphisms}

Work over the fixed category $\mathscr V$. A morphism consists of
a $\mathscr V$-enriched functor $F:\mathcal K\to\mathcal K'$
preserving restriction, the chosen binary restriction products, zeros, joins of restriction idempotents,
Turing morphism, and indexing arrow, together with bounded lattice
homomorphisms $\eta_X:P_X\to P'_{FX}$ compatible with inverse image.
Writing $F_{X,Y}:H(X,Y)\to H'(FX,FY)$ for its hom-object
morphisms, the induced ordinary functor
$F_0:\mathcal K_0\to\mathcal K'_{0}$ sends $X$ to $FX$ and
a point $\widehat f:1\to H(X,Y)$ to
$F_{X,Y}\widehat f:1\to H'(FX,FY)$.
Preservation of restriction means
\[
 r'_{FX,FY}F_{X,Y}=F_{X,X}r_{X,Y}.
\]
Compatibility of the predicate morphisms means
\[
 \eta_Xa_{X,Y}
   =a'_{FX,FY}\bigl(\mathsf T(F)_{X,Y}\times\eta_Y\bigr)
\]
on $\mathsf T(X,Y)\times P_Y$.
We require strict preservation of the chosen products, $FU=U'$,
$FS=S'$, preservation of $s,p_+,p_-$, and
\[
 \eta_Xj_X=j'_{FX}E(F)_X,
 \qquad \mathsf T(F)\rho=\rho'C(F).
\]
The arrows $E(F),\mathsf T(F),C(F)$ are induced by the equalizer and
pullback universal properties. The last equation is equivalent to
preservation of the section $\sigma$, since the inclusion of $R'$
in $C_{U'}\times\mathsf T(U',U')$ is monic.
Identity functors and composite functors, with their corresponding
predicate components, satisfy these equations.

\begin{definition}\label{def:original-definability-algebra}
A \emph{definability algebra with restriction idempotents as predicates}
has the enriched category and structure of
Section~\ref{sec:definability-algebras}, together with $s$, $\rho$, and disjoint
$a_+,a_-:1\to E_S$ satisfying
\begin{equation}\label{eq:original-algebra-representation}
 a^E_{U\otimes U,S}\langle\chi\rho,a_\pm!_{C_U}\rangle
   =\pi_2^*d_\pm.
\end{equation}
Here $a^E$ is the inverse-image action on restriction idempotents.
A morphism is an enriched functor preserving this structure,
$s,a_+,a_-$, and $\rho$.
\end{definition}

Write $\mathsf{DefAlg}(\mathscr V)$ and
$\mathsf{DefAlg}_{\mathrm{pred}}(\mathscr V)$ for the categories of
definability algebras with restriction idempotents as predicates and with general predicates, respectively.

\begin{proposition}\label{prop:algebra-comparison}
The assignment $P_X=E_X$, $j_X=1_{E_X}$, with inverse images of restriction idempotents,
defines a fully faithful functor
\[
 I:\mathsf{DefAlg}(\mathscr V)
   \longrightarrow\mathsf{DefAlg}_{\mathrm{pred}}(\mathscr V).
\]
If the distinguished predicates of a general definability algebra
factor as $p_\pm=j_Sa_\pm$ for $a_\pm:1\to E_S$, then the same
enriched structure, $s$, and $\rho$ give a definability algebra with restriction idempotents as predicates. Conversely, any compatible predicate structure $(P,j)$
extends a definability algebra with restriction idempotents as predicates by $p_\pm=j_Sa_\pm$.
\end{proposition}
\begin{proof}
The inverse-image action on restriction idempotents satisfies the required equations.
On morphisms set $I(F)=(F,E(F))$. Preservation of restriction
makes $E(F)$ compatible with inverse images. Between two objects
in the image of $I$, the equation
$\eta_Xj_X=j'_{FX}E(F)_X$ forces $\eta_X=E(F)_X$.
Every morphism between objects in the image is therefore uniquely
of the form $I(F)$, proving fullness and faithfulness.

When $p_\pm=j_Sa_\pm$, monicity of $j_S$ reflects disjointness.
Compatibility with inverse image rewrites
\eqref{eq:generalized-algebra-representation} as equalities after
composition with $j_{U\otimes U}$. Its monicity reflects those
equalities and gives~\eqref{eq:original-algebra-representation}.
Conversely, composing the latter equations with
$j_{U\otimes U}$ gives the generalized equations.
\end{proof}

The functor $\Gamma$ preserves finite limits. Consequently,
$\Gamma E_X$ is the lattice of restriction idempotents of $X$ in $\mathcal K_0$,
$\Gamma \mathsf T(X,Y)$ consists of its total arrows, and
$\Gamma C_U$ gives pairs of complementary restriction idempotents. The internal lattice
operations make $\Gamma P_X$ a bounded distributive lattice;
inverse images make these lattices a contravariant functor on
$\operatorname{Tot}(\mathcal K_0)$. The monomorphisms $j_X$
induce a natural embedding of lattices of restriction idempotents, and the section
$\sigma$ supplies a total representing arrow for each complemented restriction idempotent. In a one-object presentation with $S=U$, this gives the
one-object predicate-lattice presentation. In the case of restriction idempotents as predicates it recovers the original definability algebra and its
representing operation \cite[Section~4]{DefinabilityAlgebras}.

\subsection{Categorical inseparability}
\label{sec:categorical-representation}

Let $C_S$ be the pullback defined by
\eqref{eq:v8-complement-pullback}, with $U$ replaced by $S$,
and write $b_+,b_-:C_S\to E_S$ for its two projections.
Define $\operatorname{Sep}\rightarrowtail C_S$ as the equalizer of
\[
 \bigl\langle p_+!\wedge j_Sb_+,\ p_-!\wedge j_Sb_-\bigr\rangle,
 \quad \langle p_+!,p_-!\rangle:C_S\longrightarrow P_S\times P_S.
\]
Thus its defining equations are the separator inequalities,
expressed by meet. Define $z_U:Z_U\rightarrowtail1$ by the equalizer
\[
 Z_U\longrightarrow1
 \mathrel{\substack{\xrightarrow{\ \bot_{E_U}\ }\\[-.5ex]
                   \xrightarrow[\ \top_{E_U}\ ]{}}} E_U.
\]

\begin{theorem}[Categorical inseparability]
\label{thm:internal-inseparability}
For a definability algebra over a finitely complete category
$\mathscr V$, the structure arrow $\operatorname{Sep}\to1$
factors through $z_U:Z_U\rightarrowtail1$.
Consequently, if $\bot_{E_U}\ne\top_{E_U}$ as arrows $1\to E_U$,
there is no arrow $1\to\operatorname{Sep}$.
If $\mathscr V$ has a strict initial object and $Z_U$ is initial,
then $\operatorname{Sep}$ is initial.
\end{theorem}
\begin{proof}
Pull the structure back to $\mathscr V/\operatorname{Sep}$.
The two projections give morphisms
$a,b:\operatorname{Sep}\to E_S$ representing complementary restriction idempotents,
with $p_+\leq j_Sa$ and $p_-\leq j_Sb$.
All operations and equations below are interpreted in this slice;
their displayed parameter morphisms belong to $\mathscr V$.
Put $\delta=s\Delta_U$, where
$\Delta_U:U\to U\otimes U$ is the diagonal. Inverse image gives morphisms
\[
 d=\delta^*b,\qquad d'=\delta^*a
 \quad:\operatorname{Sep}\longrightarrow E_U
\]
representing complementary restriction idempotents on $U$, hence an arrow $\gamma:\operatorname{Sep}\to C_U$ in
$\mathscr V$. Compose it with $\rho:C_U\to \mathsf T(U,U)$ to obtain $g$.
The representation equations and separator inequalities give
\[
 \pi_2^*j_Ud\leq c_g^*j_Sa,
 \qquad
 \pi_2^*j_Ud'\leq c_g^*j_Sb,
 \qquad c_g=s(g\otimes1_U).
\]
Both pairs are complementary. In any bounded distributive
lattice object, $x\leq y$ and $x'\leq y'$ for complementary
pairs imply $x=y$: distributivity gives
$y=y\wedge(x\vee x')=(y\wedge x)\vee(y\wedge x')=x$.
This is an equation of lattice operations and therefore applies
in the slice. It follows that
$c_g^*j_Sa=\pi_2^*j_Ud$.

Pairing the constant identity with
$g=\rho\gamma:\operatorname{Sep}\to\mathsf T(U,U)$ gives
\[
 h=\langle1_U!,g\rangle:
 \operatorname{Sep}\longrightarrow\mathsf T(U,U\otimes U).
\]
Compositions of these parameterized arrows use the diagonal of
$\operatorname{Sep}$. The equations
$c_gh=\delta g$ and $\pi_2h=g$, together with compatibility of
$j$ with inverse image, yield
\[
 j_U\bigl((\delta g)^*a\bigr)
   =j_U(g^*d)
   =j_U\bigl((\delta g)^*b\bigr).
\]
Monicity of $j_U$ gives $(\delta g)^*a=(\delta g)^*b$.
These two morphisms have meet $\bot_{E_U}!_{\operatorname{Sep}}$
and join $\top_{E_U}!_{\operatorname{Sep}}$; hence
\[
 \bot_{E_U}!_{\operatorname{Sep}}
     =\top_{E_U}!_{\operatorname{Sep}}.
\]
The equalizer universal property gives the asserted factorization.
A global arrow into $\operatorname{Sep}$ would equate the two
bounds. If $Z_U$ is initial and the initial object is strict,
the resulting arrow $\operatorname{Sep}\to Z_U$ is an
isomorphism, proving the final assertion.
\end{proof}

Taking global arrows gives the usual categorical formulation.
Write $E(X)=\mathscr V(1,E_X)$ and
$P(X)=\mathscr V(1,P_X)$. These are bounded distributive lattices;
the internal actions give functors
\[
 E,P:\operatorname{Tot}(\mathcal K_0)^{\mathrm{op}}
       \longrightarrow\mathbf{BDLat},
\]
and the $j_X$ give a natural transformation $j:E\to P$ with
injective components. In this notation a total arrow
$g:U\to U$ represents a complemented restriction idempotent $e$ when
\begin{equation}\label{eq:integration-rows}
 c_g^*p_+=\pi_2^*j_U(e),\qquad
 c_g^*p_-=\pi_2^*j_U(e^c),
 \qquad c_g=s(g\otimes1_U).
\end{equation}

\begin{corollary}\label{thm:integration-diagonal}
Let $\mathcal K_0$ be a restriction category with binary restriction
products, restriction zeros, and bounded distributive lattices of restriction idempotents whose total inverse images preserve finite joins.
Let $P:\operatorname{Tot}(\mathcal K_0)^{\mathrm{op}}\to\mathbf{BDLat}$
and let $j:E\to P$ be natural with injective components.
Fix $U,S$, a total arrow $s:U\otimes U\to S$, and disjoint
$p_+,p_-\in P(S)$, with $1_U\ne0_{U,U}$.
If a complemented $D\in E(S)$ satisfies
\[
 p_+\leq j_S(D),\qquad p_-\leq j_S(D^c),
\]
then $d=(s\Delta_U)^*(D^c)$ has no representing total arrow.
In particular, representation of every complemented restriction idempotent on
$U$ excludes every such separator.
\end{corollary}
\begin{proof}
A proposed representative of $d$ gives the calculation in the
proof of Theorem~\ref{thm:internal-inseparability}, with
$a=D$ and $b=D^c$. It equates complementary restriction idempotents on $U$
and hence gives $1_U=0_{U,U}$.
\end{proof}

\begin{corollary}[Inseparability in definability algebras]
\label{cor:definability-algebra-inseparability}
If the underlying recursion category $\mathcal K_0$ of a
definability algebra has $1_U\ne0_{U,U}$, its distinguished predicates have no
separator given by a complemented restriction idempotent.
\end{corollary}
\begin{proof}
The section $\sigma$ supplies a representing total arrow for
every complemented restriction idempotent. Apply
Corollary~\ref{thm:integration-diagonal}.
\end{proof}

The distinction between the fixed objects $U$ and $S$ is
compatible with the isotype condition on a recursion category.

\begin{lemma}\label{lem:sentence-object-transport}
Let $b:S\to U$ be an isomorphism. Put
$\widetilde s=bs$ and
$\widetilde p_\pm=(b^{-1})^*p_\pm$.
These data have the same representation pullback over $C_U$,
up to its canonical isomorphism, and hence the same sections.
\end{lemma}
\begin{proof}
Isomorphisms are total. Functoriality of the internal action gives
$(bc_g)^*(b^{-1})^*p_\pm=c_g^*p_\pm$ as morphisms on $\mathsf T(U,U)$.
Thus the morphism $\ell$ defining the representation pullback
is unchanged, as is its morphism from $C_U$.
\end{proof}

\subsection{Boolean presentations}

Apply $\Gamma$ to the predicate objects to obtain the functor $P$, and
write $\mathcal A=\operatorname{Tot}(\mathcal K_0)$.
Let $\mathcal C$ be a regular extensive category with finite
coproducts, and let $J:\mathcal A\to\mathcal C$ preserve binary
products. Its subobject lattices are bounded distributive:
finite unions are images of coproducts, and regularity and
extensivity make these unions stable under inverse image.
Suppose now that $P(X)=\operatorname{Sub}_{\mathcal C}(JX)$,
with reindexing along $Jf$. A \emph{Boolean presentation} is a regular epimorphism
$q:JS\twoheadrightarrow B$, where $B$ is an internal Boolean
algebra in $\mathcal C$. Put $\mathbf2=1+1$, with its two injections
$\iota_0,\iota_1$. Its Boolean operations are determined by the
two-element truth tables. Write $\bot,\top:1\to B$ for the bounds
of $B$, and set $\iota=[\bot,\top]:\mathbf2\to B$.
Take
\[
 p_+=\operatorname{Eq}(q,\top!_{JS}),\qquad
 p_-=\operatorname{Eq}(q,\bot!_{JS}),
\]
and assume $p_+\wedge p_-=0$.

For a complemented restriction idempotent $e\in E(X)$, its image
$j_X(e)$ is complemented in $\operatorname{Sub}(JX)$. Extensivity
therefore gives a unique characteristic arrow $\chi_e:JX\to\mathbf2$
whose $1$-fiber is $j_X(e)$ and whose $0$-fiber is $j_X(e^c)$.
Indeed, the two complementary inclusions give an isomorphism
from their coproduct to $JX$. We identify $J(U\otimes U)$ with
$JU\times JU$ by the product comparison.

\begin{proposition}\label{prop:integration-family-equation}
A total arrow $g:U\to U$ represents $e$ if and only if the
following square commutes:
\begin{equation}\label{eq:integration-family-quotient}
\begin{tikzcd}[column sep=large,row sep=large]
 JU\times JU \arrow[r,"Js\,(Jg\times1)"] \arrow[d,"\chi_e\pi_2"'] & JS \arrow[d,"q"] \\
 \mathbf2 \arrow[r,"\iota"'] & B.
\end{tikzcd}
\end{equation}
\end{proposition}
\begin{proof}
The two fibers of $\chi_e\pi_2$ decompose $JU\times JU$ as a
coproduct. The representation equations say that $qJc_g$ is
$\top$ on its $1$-fiber and $\bot$ on its $0$-fiber. The
coproduct decomposition gives~\eqref{eq:integration-family-quotient}.

Conversely, that equality gives the inclusions of those two
fibers into $(Jc_g)^*p_+$ and $(Jc_g)^*p_-$, respectively.
These latter subobjects are disjoint. If complementary elements
$a,a^c$ satisfy $a\leq b$, $a^c\leq c$, and $b\wedge c=0$
in a bounded distributive lattice, then $b=a$ and $c=a^c$.
Applying this observation proves both representation equations.
\end{proof}

The representation section of a definability algebra supplies
the representation hypothesis of the following corollary.

\begin{corollary}\label{cor:integration-no-boolean-map}
Assume $1_U\ne0_{U,U}$ and every complemented restriction idempotent on $U$ is
represented. There is no arrow $h:B\to\mathbf2$ such that
$h\iota=1_{\mathbf2}$ and $hq=\chi_D$ for a complemented
$D\in E(S)$. In particular, if for every Boolean-algebra
homomorphism $h:B\to\mathbf2$ there is a complemented
$D\in E(S)$ with $hq=\chi_D$, then no such homomorphism exists.
\end{corollary}
\begin{proof}
If $hq=\chi_D$, then $h\iota=1$ gives
$p_+\leq j_S(D)$ and $p_-\leq j_S(D^c)$.
Corollary~\ref{thm:integration-diagonal} excludes this separator.
Every Boolean-algebra homomorphism preserves both bounds, so
it satisfies $h\iota=1$.
\end{proof}

\subsection{Boolean quotients and essential incompleteness}
\label{sec:boolean-extensions}

A \emph{Boolean quotient} $r:B\twoheadrightarrow B'$ is a Boolean-algebra
homomorphism whose underlying arrow is a regular epimorphism.
It induces the presentation $q_r=rq$, whose fibers over the Boolean bounds are
\[
 p^r_+=\operatorname{Eq}(q_r,\top_{B'}!_{JS}),\qquad
 p^r_-=\operatorname{Eq}(q_r,\bot_{B'}!_{JS}).
\]
The presentation $q_r$ is \emph{consistent} if
$p^r_+\wedge p^r_-=0$, and \emph{complete} if
$p^r_+\vee p^r_-=1_{\operatorname{Sub}(JS)}$.

\begin{lemma}\label{lem:complete-boolean-presentation}
For a Boolean presentation $q_r:JS\twoheadrightarrow B'$, put
$\iota_{B'}=[\bot_{B'},\top_{B'}]:\mathbf2\to B'$.
The presentation is consistent if and only if $\iota_{B'}$ is
monic, and complete if and only if $\iota_{B'}$ is a regular
epimorphism. Consequently it is consistent and complete if and
only if $\iota_{B'}$ is an isomorphism.
\end{lemma}
\begin{proof}
Each bound $1\to B'$ is monic. Their join as subobjects of $B'$
is the image of $\iota_{B'}$.
Inverse image along the regular epimorphism $q_r$ preserves finite
joins and reflects the subobject order. It follows that
$p^r_+\vee p^r_-=1$ if and only if this image is all of $B'$.
Similarly, $p^r_+\wedge p^r_-=0$ if and only if the two bound
subobjects are disjoint. In an extensive category the copairing
of two monomorphisms is monic exactly when their images are
disjoint. Thus consistency is equivalent to monicity of
$\iota_{B'}$. A monic regular epimorphism is an isomorphism.
\end{proof}

\begin{theorem}[Categorical essential incompleteness]
\label{thm:categorical-essential-incompleteness}
Let a definability algebra over $\mathscr V$ have underlying
recursion category $\mathcal K_0$ with $1_U\ne0_{U,U}$.
Let $\mathcal C$ be regular and extensive, and let
$J:\operatorname{Tot}(\mathcal K_0)\to\mathcal C$ preserve binary
products. Suppose that the global predicate functor is naturally
identified with $\operatorname{Sub}_{\mathcal C}(J{-})$.
Let $q:JS\twoheadrightarrow B$ be a Boolean presentation whose
fibers over the Boolean bounds are the distinguished predicates $p_+,p_-$.

Let $r:B\twoheadrightarrow B'$ be a Boolean quotient whose
induced presentation $q_r=rq$ is consistent. If there are
restriction idempotents $a_+,a_-:S\to S$ satisfying
\[
 p^r_+=j_S(a_+),\qquad p^r_-=j_S(a_-)
\]
then $q_r$ is incomplete. Equivalently,
$[\bot_{B'},\top_{B'}]:\mathbf2\to B'$ is not a regular
epimorphism.
\end{theorem}
\begin{proof}
Suppose that the presentation is also complete. Since $j_S$ is an
injective bounded lattice homomorphism, consistency and
completeness give
\[
 a_+a_-=0_{S,S},\qquad a_+\vee a_-=1_S.
\]
Thus $a_-$ is the complement of $a_+$. Preservation of the Boolean
bounds by $r$ gives
\[
 p_+\leq p^r_+=j_S(a_+),\qquad
 p_-\leq p^r_-=j_S(a_-).
\]
The complemented restriction idempotent $a_+$ therefore separates $p_+$ and
$p_-$. The categorical inseparability theorem excludes this
separator. The last assertion follows from
Lemma~\ref{lem:complete-boolean-presentation}.
\end{proof}

\begin{proposition}\label{prop:complete-extensions-and-decisions}
For a Boolean presentation $q:JS\twoheadrightarrow B$, the
following conditions are equivalent.
\begin{enumerate}
\item There is a Boolean quotient $r:B\twoheadrightarrow B'$
such that the induced presentation $rq$ is consistent and complete,
with both fibers over the Boolean bounds in the image of $j_S$.
\item There are a Boolean-algebra homomorphism
$h:B\to\mathbf2$ and a complemented restriction idempotent $D$ on
$S$ such that $hq=\chi_D$.
\end{enumerate}
\end{proposition}
\begin{proof}
For (1), Lemma~\ref{lem:complete-boolean-presentation} makes
$\iota_{B'}$ an isomorphism. It is a Boolean-algebra homomorphism,
so $h=\iota_{B'}^{-1}r$ is also one. If the fibers over the Boolean bounds are
$j_S(a_+)$ and $j_S(a_-)$, injectivity of $j_S$ makes $a_+$ and
$a_-$ complementary. The two fibers of $hq$ are these subobjects,
so $hq=\chi_{a_+}$.

Conversely, every Boolean-algebra homomorphism
$h:B\to\mathbf2$ satisfies $h[\bot_B,\top_B]=1_{\mathbf2}$.
It is therefore a split epimorphism and hence a regular
epimorphism. Its presentation $hq=\chi_D$ is consistent and
complete, with fibers over the Boolean bounds $j_S(D)$ and $j_S(D^c)$.
\end{proof}

In particular, the hypotheses of
Theorem~\ref{thm:categorical-essential-incompleteness} exclude every
Boolean-algebra homomorphism $h:B\to\mathbf2$ whose composite
with $q$ is the characteristic arrow of a complemented
restriction idempotent. If every Boolean-algebra homomorphism to
$\mathbf2$ has this property, there is no such homomorphism.

\section{The numerical interpretation}
\label{sec:integration-numerical}

Let $\mathcal R_{\mathrm{num}}$ have the infinite recursive subsets
of $\mathbb N$ as objects and partial recursive maps as arrows,
with fixed recursive tuple coding. Universal partial recursive
application and the parameter theorem
\cite[Section~3.4.1]{AspertiLongo} make this a recursion
category \cite[Section~10(1)]{Heller}. Every object is recursively
isomorphic to $\mathbb N$, and binary products remain in this
category. Let $\mathcal A$ be its total-map category, and let
$E(X)$ be the lattice of restriction idempotents on $X$ in
$\mathcal R_{\mathrm{num}}$.
Each such idempotent is the partial identity $e_A$ on a c.e.
subset $A\subseteq X$. We identify $e_A$ with $A$ in this section;
under this identification, $E(X)$ is the bounded distributive
lattice of c.e. subsets of $X$, and inverse image is recursive
inverse image.

A \emph{partial equivalence relation} (PER) is a symmetric and
transitive binary relation. We use the PER category of
\cite[Definition~3.4.1.1]{AspertiLongo}. Let $\mathcal P$ have all PERs on
$\mathbb N$ as objects, with arrows defined by partial recursive
trackers as follows. An arrow
$R\to S$ is represented by a partial recursive function $f$ defined
on $|R|=\{n:nRn\}$ such that $nRm$ implies $f(n)S f(m)$.
Two trackers define the same arrow if their values are $S$-related
at every supported input. Neither a relation nor its support is
required to be c.e. For $A\subseteq\mathbb N$, write $D_A$ for
equality restricted to $A$; tuple coding gives the corresponding
notation for subsets of numerical objects. The functor
$J:\mathcal A\to\mathcal P$ sends an object to its equality
relation and a recursive map to its tracker. It is faithful and
preserves binary products; on these objects it is full.
A partial recursive map $f:X\rightharpoonup Y$ is represented
by the span $D_X\leftarrow D_{\operatorname{dom}(f)}\to D_Y$.
Its left leg is a c.e. domain inclusion. Equality and
composition of these spans reproduce equality and composition of
partial recursive maps.

To compare $\mathcal P$ with the modest-set presentation, send $R$
to its set of equivalence classes $|R|/R$, with a class realized
by its members, and send a tracker $f$ to $[n]\mapsto[f(n)]$.
Conversely, the disjoint nonempty realizing sets of a modest set
define a PER on their union; the same trackers define its arrows.
These assignments give the equivalence with modest sets over
Kleene's first partial combinatory algebra
\cite[Proposition~1.2.2]{Bauer}. It transfers the finite-limit,
regular-image, and coproduct constructions of
\cite[Sections~1.3.1--1.3.3]{Bauer} to $\mathcal P$.
Thus $\mathcal P$ is regular and extensive.
Consequently $P(X)=\operatorname{Sub}_{\mathcal P}(JX)$, with
pullback along $Jf$, defines a functor to bounded distributive
lattices.

\begin{lemma}\label{lem:integration-numerical-subobjects}
There is a natural transformation $j:E\Rightarrow P$ whose
components are injective bounded lattice homomorphisms. Its
component at $X$ sends a c.e. subset $A\subseteq X$ to the
subobject represented by $i_A:D_A\hookrightarrow D_X$.
More generally, arbitrary support
inclusions satisfy
\[
 [i_A]\leq[i_B]\ \Longleftrightarrow\ A\subseteq B,
 \qquad
 (Jf)^*[i_A]=[i_{f^{-1}A}].
\]
The arrows $D_{\mathbb N}\to\mathbf2=D_{\{0,1\}}$ are exactly
total recursive Boolean functions, and their fibers over $1$ are exactly the
images under $j$ of recursive subsets. These are also exactly
the complemented subobjects of $D_{\mathbb N}$.
\end{lemma}
\begin{proof}
A factorization of $i_A$ through $i_B$ must send each $a$ to $a$.
It exists precisely when $A\subseteq B$, and is then tracked by
the identity. For a pullback, restrict the source equality
relation to $f^{-1}A$; the defining square and its universal
property use the same trackers as the corresponding maps into
$D_X$ and $D_A$. This proves the displayed assertions, naturality,
and preservation of meets and bounds.

For c.e. sets $A,B$, their categorical join is the image of
$D_A+D_B\to D_X$. An input in $A\cup B$ can be assigned to a
summand by dovetailing enumerations of $A$ and $B$ until that
input appears. This gives a recursive tracker on $A\cup B$ and
identifies the image with $D_{A\cup B}$. Hence $j$ preserves
joins. Finally, a tracker from $D_{\mathbb N}$ to $\mathbf2$
is a total recursive Boolean function; its pullback of $1$ is
the indicated recursive support inclusion. Conversely every
recursive subset has such a characteristic function. A complemented
subobject and its complement give a coproduct decomposition of
$D_{\mathbb N}$; its inverse computes the summand tag and hence
a Boolean function with that subobject as its fiber over $1$. This proves
the final assertion.
\end{proof}

The restriction to c.e. sets in the join assertion matters.
For a nonrecursive $A\subseteq\mathbb N$, the mono
$D_A+D_{\mathbb N\setminus A}\to D_{\mathbb N}$ is the join
of its two summands but is not invertible: an inverse would
compute the summand tag and decide $A$. Thus arbitrary support
inclusions do not identify the full subobject lattice with a
powerset.

Let $\mathcal E_{\mathrm{rec}}$ be the full subcategory of $\mathcal P$
on c.e. PERs. For $R,S$ in this subcategory, a tracker $f:R\to S$
determines the c.e. relation
\[
 F_f(x,y)\quad\Longleftrightarrow\quad
 xRx,\quad f(x)\downarrow,\quad f(x)S y.
\]
It is saturated under $R$ and $S$, total on $|R|$, and single-valued
modulo $S$. Conversely, searching a c.e. relation
$F\subseteq |R|\times |S|$ with these properties for an output
gives a partial recursive tracker.
The assignments are inverse on arrows and preserve composition.
They identify $\mathcal E_{\mathrm{rec}}$ with Lambek and Scott's
category of c.e. PERs and c.e. functional relations
\cite[Definition~2.4 and Remark~2.6]{LambekScott}, which is exact
\cite[Theorem~4.5]{LambekScott}.
The finite-limit, coproduct, and regular-image constructions preserve
c.e. relations, so its inclusion in $\mathcal P$ preserves these
constructions and $\mathcal E_{\mathrm{rec}}$ is extensive.

The connection with completion of the numerical recursion category
can be stated explicitly. Let $\mathcal D_{\mathrm{rec}}$ have c.e.
subsets $A\subseteq\mathbb N$ as objects and partial recursive
functions total on $A$ as arrows $A\to B$, identified by equality
on $A$. Write $\operatorname{Split}_r(\mathcal R_{\mathrm{num}})$
for the category obtained by splitting restriction idempotents:
an object is such an idempotent $e:X\to X$, and an arrow from $e$
to $d$ is an arrow $f:X\to Y$ with $f=df=fe$.
Its identity at $e$ is $e$, and totality there means $\bar f=e$.

\begin{proposition}\label{prop:numerical-exact-completion}
There are equivalences
\[
 \operatorname{Tot}\bigl(\operatorname{Split}_r
       (\mathcal R_{\mathrm{num}})\bigr)
       \simeq\mathcal D_{\mathrm{rec}},
 \qquad
 \operatorname{Ex}_{\mathrm{reg}}(\mathcal D_{\mathrm{rec}})
       \simeq\mathcal E_{\mathrm{rec}},
\]
where $\operatorname{Ex}_{\mathrm{reg}}$ denotes exact completion
of a regular category.
\end{proposition}
\begin{proof}
A restriction idempotent is partial identity on a c.e. set $A$,
and every c.e. set occurs as such a domain in $\mathbb N$.
An arrow total relative to this idempotent has domain exactly $A$.
Conversely, a tracker total on $A$ can be restricted to precisely
$A$ by first waiting for its input to appear in an enumeration
of $A$. This gives the first equivalence.

Finite limits in $\mathcal D_{\mathrm{rec}}$ have c.e. supports.
For $f:A\to B$, the image $f(A)$ is c.e., and the surjection
$A\to f(A)$ has a partial recursive section: enumerate $A$ and
dovetail the computations until a preimage of the input is found.
Surjections are therefore regular epimorphisms. Conversely,
factoring a regular epimorphism through its image shows that it
is surjective. Surjections are stable under pullback, so
$\mathcal D_{\mathrm{rec}}$ is regular.

Every mono is injective, as tested by maps from singleton supports,
and is isomorphic to the inclusion of its c.e. image; the inverse on that image is obtained by the
same search. Thus its internal equivalence relations are
precisely c.e. equivalence relations on c.e. supports.
The relation construction of exact completion
\cite{CarboniVitale} has these as objects and saturated total
single-valued relations as arrows. Its objects are therefore
c.e. PERs, and the tracker comparison above identifies its
arrows with those of $\mathcal E_{\mathrm{rec}}$.
\end{proof}

\subsection{Strong numeralwise definability}

Fix a consistent theory $T$ in a classical first-order language
with recursive syntax, a recursive sequence of closed numeral terms
$\nu_n$, and recursive substitution. Theories are identified
with their deductive closures, and all extensions below are in
this language. Let $\mathsf{Sent}\subseteq\mathbb N$ be the infinite
recursive set of G\"odel codes of sentences, and fix a recursive
bijection $\beta:\mathbb N\to\mathsf{Sent}$. Write $\sigma_b$
for the sentence with G\"odel code $\beta(b)$. Take $S=U=\mathbb N$.
Decode $a\in\mathbb N$ as a formula $\varphi_a(x)$
whose only possible free variable is $x$; invalid codes denote a
fixed contradiction. Numeral substitution defines the total
recursive arrow $s:U\otimes U\to U$ by
$s(a,n)=\beta^{-1}(\ulcorner\varphi_a(\nu_n)\urcorner)$.
Thus every element of $U$ indexes a sentence through $\beta$,
while the first argument of $s$ is decoded as a unary-formula code.
Let $p_T^+,p_T^-\in P(S)$ be the support inclusions of
\begin{align*}
 \mathsf{Th}_T&=\{b\in\mathbb N:T\vdash\sigma_b\},\\
 \mathsf{Rf}_T&=\{b\in\mathbb N:T\vdash\neg\sigma_b\}.
\end{align*}
They are disjoint by consistency, although they need not lie in
the image of $j$. For a total recursive $g:\mathbb N\to\mathbb N$,
write $c_g(x,n)=s(g(x),n)$.

\begin{proposition}\label{prop:integration-arbitrary-adequacy}
For a recursive set $A\subseteq\mathbb N$, the following are
equivalent.
\begin{enumerate}
\item There is a total recursive $g$ satisfying the representation equations
\[
 c_g^*p_T^+=\pi_2^*j(A),\qquad
 c_g^*p_T^-=\pi_2^*j(\mathbb N\setminus A).
\]
\item Some formula $\varphi_a(x)$ strongly defines $A$
numeralwise in $T$: each $n\in A$ satisfies
$T\vdash\varphi_a(\nu_n)$, and each $n\notin A$ satisfies
$T\vdash\neg\varphi_a(\nu_n)$.
\end{enumerate}
Consequently, representation of every complemented restriction idempotent is exactly strong numeralwise definability
of every recursive set, without a c.e. hypothesis on $T$.
\end{proposition}
\begin{proof}
All four predicates in the representation equations are support inclusions.
Lemma~\ref{lem:integration-numerical-subobjects} identifies their
equality with equality of their supports. Thus the first equation says, for every $x,n$,
$T\vdash\varphi_{g(x)}(\nu_n)\quad\Longleftrightarrow\quad n\in A$,
and the second says the analogous equivalence for refutation and
$n\notin A$. Taking $x=0$ yields (2). Conversely, if $a$
strongly defines $A$, take the constant recursive function
$g(x)=a$. Consistency excludes a proof of the opposite instance
in each case, so the defining implications are equivalences and
give both representation equations.
\end{proof}

\begin{corollary}[Numerical realization of the representation section]
\label{cor:numerical-definability-algebras}
Let $\mathcal K$ be the $\mathbf{Set}$-enriched category with
hom-objects $H(X,Y)=\mathcal R_{\mathrm{num}}(X,Y)$. Evaluation at
the unique element of $1$ identifies $\mathcal K_0$ with
$\mathcal R_{\mathrm{num}}$. Equip $\mathcal K$ with
the bounded distributive lattice objects
$P_X=\operatorname{Sub}_{\mathcal P}(JX)$ in $\mathbf{Set}$ and
the embedding of restriction idempotents $j$ of
Lemma~\ref{lem:integration-numerical-subobjects}.
Fix a consistent theory $T$ with the effective syntax and numeral
terms above, and take $S=U=\mathbb N$ and
$s,p_T^+,p_T^-$ as defined above.
The projection $q_R:R\to C_U$ from the representation pullback
has a section if and only if every recursive subset of
$\mathbb N$ is strongly definable numeralwise in $T$.
When these conditions hold and $T$ is c.e., the same chosen
representing arrows give a definability algebra with predicate
functor $E$ and distinguished restriction idempotents
$e_{\mathsf{Th}_T}$ and $e_{\mathsf{Rf}_T}$.
\end{corollary}
\begin{proof}
Composition, restriction, and pairing are maps between the
indicated hom-objects in $\mathbf{Set}$.
Fix an acceptable program numbering and its parameter function
$k$. For each partial recursive arrow
$f:\mathbb N\times\mathbb N\rightharpoonup\mathbb N$, let
$e_f$ be its least program index. The total recursive arrow
$\Lambda(f)(x)=k(e_f,x)$ satisfies
$\tau(\Lambda(f)\times1)=f$. This defines the required morphism
$\Lambda:H(U\otimes U,U)\to\mathsf T(U,U)$ in $\mathbf{Set}$.
The total-arrow action on predicates is the map
\[
 \mathsf T(X,Y)\times P(Y)\longrightarrow P(X),
 \qquad (f,a)\longmapsto (Jf)^*a.
\]
Functoriality of inverse image and
Lemma~\ref{lem:integration-numerical-subobjects} give its action
laws and compatibility with $j$. Joins of restriction idempotents
correspond to unions of their domains, which are stable under
arbitrary partial recursive inverse images.
Thus the equalizers and pullbacks defining $C_U$ and $R$ apply
to these data.

The two projections from $C_U$ give partial identities on
complementary c.e. supports $A$ and $\mathbb N\setminus A$.
Such an $A$ is recursive, and every recursive $A$ occurs.
Suppose that every recursive set is strongly definable. Let
$a_A$ be the least code of a formula strongly defining $A$ in
$T$, and let $g_A:\mathbb N\to\mathbb N$ be the constant
recursive function with value $a_A$.
Proposition~\ref{prop:integration-arbitrary-adequacy} verifies
the two representation equations. Therefore the assignment
\[
 (e_A,e_{A^c})\longmapsto
 \bigl((e_A,e_{A^c}),g_A\bigr)
\]
defines a morphism $\sigma:C_U\to R$ in $\mathbf{Set}$.
Its composite with $q_R$ is the identity, so it is a section.
The representing operation is the composite of $\sigma$ with
the projection $R\to \mathsf T(U,U)$.

Conversely, a section supplies a representing arrow $g_A$ for
each recursive $A$. Evaluating $g_A$ at $0$ gives the code of
a formula strongly defining $A$, by
Proposition~\ref{prop:integration-arbitrary-adequacy}.

If $T$ is c.e., both distinguished supports are c.e. and define
restriction idempotents on $S$. Their images under $j_S$ are
$p_T^+$ and $p_T^-$. Compatibility of $j$ with inverse image
and its injectivity identify the representation equations with
the corresponding equations in $E$. Hence the same representing
arrows give the definability algebra with predicate functor $E$.
\end{proof}


Two disjoint subsets of $\mathbb N$ are \emph{recursively inseparable}
when no recursive set contains the first and is disjoint from the second.
If every recursive set is represented, Corollary~\ref{thm:integration-diagonal}
therefore implies recursive inseparability of
$\mathsf{Th}_T$ and $\mathsf{Rf}_T$.

\subsection{The Lindenbaum quotient and essential incompleteness}

Define the PER on $S=U=\mathbb N$ by
\[
 bR_Tc\quad\Longleftrightarrow\quad
 T\vdash\sigma_b\leftrightarrow\sigma_c.
\]
The identity tracker gives the regular quotient
$q_T:JU\to B_T=(\mathbb N,R_T)$.
Indeed, its kernel pair is the equality PER on
$\{(b,c):bR_Tc\}$ with its two projections. A tracker coequalizing
those projections respects $R_T$ and descends uniquely to $B_T$,
so $q_T$ is their coequalizer. Negation is represented by the
recursive function $b\mapsto\beta^{-1}(\ulcorner\neg\sigma_b\urcorner)$;
conjunction and the bounds are encoded through $\beta$ in the same way.
These operations make $B_T$ a Boolean algebra in $\mathcal P$. Its pullbacks along
top and bottom are the theorem and refutation support inclusions.
Thus they are precisely $p_T^+,p_T^-$.
For a recursive characteristic map $d:D_{\mathbb N}\to\mathbf2$
and a formula code $a$, write $s_a:D_{\mathbb N}\to JS$ for the
arrow tracked by $n\mapsto s(a,n)$.
The formula $\varphi_a$ strongly defines $d^{-1}(1)$ numeralwise
in $T$ if and only if
\[
 q_Ts_a=[\bot,\top]d.
\]
Equality here requires a $T$-proof for each numeral instance; it
does not assert a single universally quantified $T$-theorem.

A theory is \emph{decidable} when its theorem set is recursive,
and \emph{complete} when it proves either $\sigma$ or $\neg\sigma$
for every sentence $\sigma$. It is \emph{essentially undecidable}
when every consistent extension is undecidable, and
\emph{essentially incomplete} when every consistent c.e. axiomatized
extension is incomplete.

Boolean homomorphisms $B_T\to\mathbf2$ are in bijection with
complete decidable consistent extensions of $T$. Indeed, the
composite with $q_T$ is a recursive Boolean function on
sentences. Its inverse image of $1$ contains $T$, decides each sentence,
and is closed under deduction: any proof from finitely many
sentences in this inverse image gives a $T$-provable implication from their
conjunction, and a Boolean homomorphism preserves this order.
Conversely the recursive characteristic function of a complete
decidable consistent extension respects $R_T$ and the Boolean
operations, so is a tracker for the required homomorphism.

\begin{theorem}[Numerical realization of essential incompleteness]
\label{thm:integration-essential-equivalence}
Let $T$ be a consistent classical first-order theory with the
effective syntax and numeral terms specified above. Its
Lindenbaum presentation $q_T:JS\to B_T$ in $\mathcal P$ has the
following properties.
\begin{enumerate}[label=\textup{(\roman*)}]
\item Every complemented restriction idempotent on $\mathbb N$ is
represented relative to $s,p_T^+,p_T^-$ if and only if every
recursive subset of $\mathbb N$ is strongly definable numeralwise
in $T$. These conditions are also equivalent to the existence of a
section of the representation projection $q_R:R\to C_U$
for the numerical data of
Corollary~\ref{cor:numerical-definability-algebras}.
\item There is no complemented restriction idempotent on $S$
separating $p_T^+$ and $p_T^-$ if and only if
$\mathsf{Th}_T$ and $\mathsf{Rf}_T$ are recursively inseparable.
\item There is no Boolean homomorphism $B_T\to\mathbf2$ in
$\mathcal P$ if and only if $T$ is essentially undecidable,
equivalently essentially incomplete.
\end{enumerate}
Under the equivalent conditions in \textup{(i)}, the conditions
in \textup{(ii)} and \textup{(iii)} hold.
\end{theorem}
\begin{proof}
Part (i) follows from Proposition~\ref{prop:integration-arbitrary-adequacy}
and Corollary~\ref{cor:numerical-definability-algebras}.
For (ii), a complemented restriction idempotent on $S=U$ is the
partial identity on a recursive subset of $\mathbb N$. The separator
inequalities are exactly inclusion of $\mathsf{Th}_T$ and
disjointness from $\mathsf{Rf}_T$, giving both directions.

For (iii), the preceding correspondence identifies Boolean
homomorphisms with complete decidable consistent extensions.
A complete consistent c.e. axiomatized theory is decidable, by
searching for a proof of a sentence or its negation; a decidable
theory is c.e. axiomatized by its theorems. Every consistent
decidable theory has a complete decidable consistent extension
\cite[Section~2.7, Exercise~2.50]{Mendelson}.
Applying this fact to a decidable extension of $T$ proves the
equivalence with essential undecidability and essential incompleteness.

Under (i), Corollary~\ref{thm:integration-diagonal} excludes the
separators in (ii). Every arrow $JS\to\mathbf2$ is the
characteristic arrow of a recursive subset of $S$, by the same
tracker argument as in
Lemma~\ref{lem:integration-numerical-subobjects}. Therefore
Corollary~\ref{cor:integration-no-boolean-map} excludes the Boolean
homomorphisms in (iii).
\end{proof}

For c.e. $T$, the PER $R_T$, its sentence presentation, and its
two fibers belong to $\mathcal E_{\mathrm{rec}}$. Fullness of its
inclusion in $\mathcal P$ gives the same Boolean-homomorphism
characterization in that category.

\begin{proposition}[Lindenbaum quotients induced by theory extensions]
\label{prop:numerical-boolean-extensions}
Let $T$ be a consistent classical first-order theory with the
effective syntax and numeral substitution fixed above.
For every consistent extension $T'$ of $T$ in the same language,
the identity on sentence codes induces a regular-epimorphic
Boolean-algebra homomorphism
\[
 r_{T,T'}:B_T\twoheadrightarrow B_{T'},\qquad
 r_{T,T'}q_T=q_{T'}.
\]
The fibers of $q_{T'}$ over the Boolean bounds lie in the image of $j_S$ if and
only if $T'$ is c.e. axiomatizable. Under these equivalent
conditions, $q_{T'}$ is complete if and only if $T'$ is complete.

There is a Boolean quotient $r:B_T\twoheadrightarrow B$ for
which $rq_T$ is consistent and complete and its fibers over the
bounds lie in the image of $j_S$ if and only if $T$ has a
complete consistent c.e. extension. Consequently the absence of
such quotient presentations is equivalent to essential incompleteness
of $T$.
\end{proposition}
\begin{proof}
An equivalence provable in $T$ is provable in $T'$, so the
identity on sentence codes respects the two Lindenbaum PERs.
It induces $r_{T,T'}$, preserves the Boolean operations, and
satisfies $r_{T,T'}q_T=q_{T'}$. Since $q_{T'}$ is a regular
epimorphism, factoring $r_{T,T'}$ into its regular image shows
that $r_{T,T'}$ is a regular epimorphism.

The fibers over the Boolean bounds have supports $\mathsf{Th}_{T'}$ and
$\mathsf{Rf}_{T'}$. They lie in the image of $j_S$ exactly when
these supports are c.e. A c.e. axiomatization enumerates both
supports. Conversely, a c.e. theorem set itself gives a c.e.
axiomatization. Consistency makes the two supports disjoint.
For c.e. $T'$, preservation of joins by $j_S$ identifies
$p^{T'}_+\vee p^{T'}_-=1$ with
$\mathsf{Th}_{T'}\cup\mathsf{Rf}_{T'}=\mathbb N$,
which is completeness of $T'$. Equivalently, when $T'$ is
complete, simultaneous proof searches for a sentence and its
negation give the recursive characteristic arrow of its theorem
set and the inverse to $[\bot,\top]:\mathbf2\to B_{T'}$.

A complete consistent c.e. extension therefore gives a Boolean quotient with the stated properties. Conversely, by
Proposition~\ref{prop:complete-extensions-and-decisions}, such a
Boolean quotient gives $h:B_T\to\mathbf2$ with $hq_T$ a
recursive characteristic arrow on sentences. Its $1$-fiber
contains the theorems of $T$, decides every sentence, and is
closed under deduction: a proof from finitely many sentences
gives a $T$-provable implication from their conjunction, and
$h$ preserves conjunction and Boolean order. Since $h$ preserves
both bounds, this theory is consistent. It is decidable and
hence c.e. axiomatizable. This proves both implications.
\end{proof}

\subsection{Extensions of arithmetic}
\label{sec:arithmetic-representation}

Let PA denote Peano arithmetic. In this subsection take
$\nu_n=\overline n$, the standard arithmetic numeral for $n$.
Every recursive set is strongly definable numeralwise in PA,
by representability of recursive characteristic functions
\cite[Section~16.2]{BBJ}.

\begin{corollary}\label{cor:arithmetic-representation-incompleteness}
Let $T$ be a consistent extension of PA in an effective classical
language. Every complemented restriction idempotent on $\mathbb N$
is represented relative to $s,p_T^+,p_T^-$ by a constant total arrow.
The sets $\mathsf{Th}_T$ and $\mathsf{Rf}_T$ are recursively inseparable,
and there is no Boolean homomorphism $B_T\to\mathbf2$ in $\mathcal P$.
Consequently $T$ is essentially undecidable and essentially incomplete.
The numerical representation projection admits a section, and
for c.e. $T$ the same representing arrows give a definability
algebra with predicate objects $E_X$.
\end{corollary}
\begin{proof}
Numeralwise definitions in PA remain definitions in $T$.
Proposition~\ref{prop:integration-arbitrary-adequacy} supplies
the constant representing arrows, and
Theorem~\ref{thm:integration-essential-equivalence} gives the incompleteness
conclusions. Corollary~\ref{cor:numerical-definability-algebras} supplies
the algebraic structures.
\end{proof}

\begin{corollary}\label{cor:proper-predicate-extension}
There is a definability algebra over $\mathbf{Set}$ with underlying
recursion category $\mathcal R_{\mathrm{num}}$ whose distinguished predicates lie
outside the image of the embedding of restriction idempotents. Consequently the
functor $I$ of Proposition~\ref{prop:algebra-comparison} for $\mathscr V=\mathbf{Set}$ is not essentially surjective.
\end{corollary}
\begin{proof}
Take $T=\operatorname{Th}(\mathbb N)$, the theory of the standard
model of arithmetic. It is a complete consistent extension of PA,
so Corollary~\ref{cor:arithmetic-representation-incompleteness}
and Corollary~\ref{cor:numerical-definability-algebras} supply
the representation section. Each of the theorem and refutation
supports is the inverse image of the other under recursive
negation. If either support were c.e., both would be c.e.;
completeness would then make the theorem set recursive,
contrary to their recursive inseparability.
Lemma~\ref{lem:integration-numerical-subobjects} identifies the
image of $j$ with the c.e. support inclusions. Neither
distinguished predicate belongs to it. In the image of the
functor $I$, both distinguished predicates belong to the image
of $j$. Compatibility of a definability-algebra isomorphism with
the embeddings of restriction idempotents preserves this property, proving the claim.
\end{proof}


\section{An equivariant recursion category}
\label{sec:equivariant-model}

Di Paola and Heller construct recursion categories of free actions
of finite groups \cite[Section~7.9(3), Theorem~7.9.9]{DPH}.
We use the one-object form of their construction over the numerical
recursion category and give its Boolean presentation.
Let $G$ be a finite nontrivial group, with fixed
effective multiplication and coding, and let
\[
 U=G\times\mathbb N,
 \qquad k\cdot(g,n)=(kg,n).
\]
Consider the one-object category $\mathcal R_G$ whose arrows are
the $G$-equivariant partial recursive maps $U\rightharpoonup U$.
Equivariance includes invariance of the domain.  Restriction is
partial identity on the domain, and zero is the nowhere-defined
map.  Every domain has the form $G\times D$ for a c.e. set
$D\subseteq\mathbb N$.  Conversely each such set is a domain.
Consequently the lattice of restriction idempotents is the lattice of c.e. subsets
of $\mathbb N$, with joins given by union and stable under
partial inverse image.

In Di Paola and Heller's construction, the free objects are
$G\times X$ for $X$ in $\mathcal R_{\mathrm{num}}$.
A recursive bijection $h_X:\mathbb N\to X$ gives the equivariant
isomorphism $(g,n)\mapsto(g,h_X(n))$ from $U$ to $G\times X$.
Their free-action category is therefore equivalent to its full
one-object subcategory on $U$, whose arrows are exactly those of
$\mathcal R_G$.

The product $U\times U$ carries the diagonal action.  Choose a
computable bijection
$b:\mathbb N\times G\times\mathbb N\to\mathbb N$.
Then
\[
 ((g,m),(h,n))\longmapsto(g,b(m,g^{-1}h,n))
\]
is an equivariant computable bijection $U\times U\cong U$.
Transporting the coordinatewise product of partial maps along this
bijection gives binary restriction products on $\mathcal R_G$,
written $\otimes$. These are the near products in the dominical
terminology. In the formulas below, we use the unencoded pairs
in $U\times U$.

\begin{proposition}\label{prop:equivariant-recursion}
The category $\mathcal R_G$ is a dominical recursion category
with $1_U\ne0_{U,U}$. Its lattices of restriction idempotents are bounded distributive, and finite joins
are stable under inverse image along total arrows.
\end{proposition}
\begin{proof}
Let $(\psi_a)$ be an acceptable enumeration of partial recursive
maps $G\times\mathbb N\rightharpoonup G\times\mathbb N$.
Define equivariant application by
\[
 \tau((g,a),(h,n))=(gk,m)
 \quad\text{when}\quad
 \psi_a(g^{-1}h,n)=(k,m).
\]
Let $f:U\times U\rightharpoonup U$ be equivariant and partial
recursive.  There are partial recursive functions $k_f,m_f$,
with the same domain, such that
\[
 f((g,r),(h,n))
   =\bigl(gk_f(r,g^{-1}h,n),m_f(r,g^{-1}h,n)\bigr).
\]
Fix a parameter function for the acceptable numbering, and
choose the least program index for the partial recursive map
$(r,d,n)\mapsto(k_f(r,d,n),m_f(r,d,n))$.
The parameter theorem \cite[Section~3.4.1]{AspertiLongo}
then specifies a total recursive function $a_f$ such that
\[
 \psi_{a_f(r)}(d,n)=(k_f(r,d,n),m_f(r,d,n)).
\]
Thus $\lambda_f(g,r)=(g,a_f(r))$ is total and equivariant, and
$\tau(\lambda_f\times1)=f$. Thus $\tau$ is a Turing morphism.
The assignment $\Lambda(f)=\lambda_f$ is a morphism of the
indicated hom-objects in $\mathbf{Set}$ and supplies the
indexing structure of Section~\ref{sec:definability-algebras}.

The partial-map restriction and product axioms hold directly.
Products of nonzero arrows are nonzero because both domains
are nonempty.  If an arrow $f$ is not total, choose an orbit
$G\times\{n\}$ outside its domain.  The total equivariant map
$(g,r)\mapsto(g,n)$ is nonzero and is annihilated by $f$.
Conversely a total partial map annihilates no nonzero partial
map.  These facts give the additional dominical conditions.
Finally the identity and nowhere-defined map are different.
\end{proof}

\subsection{The Boolean quotient}

Let $\mathcal K$ be the $\mathbf{Set}$-enriched category with
hom-object $H(U,U)=\mathcal R_G(U,U)$ and the composition and
restriction operations just defined. Evaluation at the unique
element of $1$ identifies $\mathcal K_0$ with $\mathcal R_G$.
Use the regular extensive category $\mathcal E_{\mathrm{rec}}$
of c.e. PERs defined in Section~\ref{sec:integration-numerical}.
Define $\mathcal C=\mathcal E_{\mathrm{rec}}^{G}$,
the category of objects of $\mathcal E_{\mathrm{rec}}$ equipped
with a $G$-action and equivariant arrows. Finite limits, regular
image factorizations, and finite coproducts are computed on the
underlying objects with their induced actions. For a regular
image factorization, the action descends uniquely to the image;
uniqueness gives the group laws. Kernel-pair coequalizers and
their pullbacks are computed on the underlying objects as well.
Hence $\mathcal C$ is regular and extensive.

The free action on $U$ defines an object of $\mathcal C$.
Total arrows of $\mathcal R_G$ embed fully faithfully as its
equivariant total recursive endomorphisms.  The corresponding
functor $J:\operatorname{Tot}(\mathcal R_G)\to\mathcal C$
preserves binary products. Its value $JU$ is the equality PER on
$G\times\mathbb N$ with its free action. Put
$P(X)=\operatorname{Sub}_{\mathcal C}(JX)$, with inverse image
along $Jf$. For a c.e. subset $D\subseteq\mathbb N$, let
$e_{G\times D}$ be the partial identity on $G\times D$, and define
$j(e_{G\times D})$ to be the inclusion
$D_{G\times D}\hookrightarrow JU$.
These inclusions preserve and reflect the order on restriction
idempotents and preserve meets, joins, bounds, and total inverse
image.  For unions, dovetail enumerations of $D$ and $E$ using
only the orbit index $n$ to select a coproduct tag.  This gives
an equivariant recursive witness for the union
$G\times(D\cup E)$; the choice of tag is independent of the
group coordinate. Thus they give the lattice embeddings $j$ required
in Section~\ref{sec:integration}.

Fix a consistent c.e. theory $T$ with effective syntax and
numeral substitution, and suppose every recursive subset of
$\mathbb N$ is strongly numeralwise definable in $T$.
Corollary~\ref{cor:arithmetic-representation-incompleteness} supplies this
hypothesis for every consistent c.e. extension of PA.
Let $B_T$ be its Lindenbaum Boolean algebra from
Section~\ref{sec:integration-numerical}, now with the trivial
$G$-action. Its bounds are equivariant global arrows.
For this example take $S=U$, and define
\[
 q_T^G:JS\longrightarrow B_T,
 \qquad q_T^G(g,a)=[a]_T.
\]
Here $[a]_T$ is the class of $a$ in the PER $R_T$;
its sentence is $\sigma_a$, with G\"odel code $\beta(a)$.
The underlying arrow is the composite of the projection
$G\times\mathbb N\to\mathbb N$ and the quotient
$q_T:D_{\mathbb N}\to B_T$.
The projection has a recursive section, and $q_T$ is a
regular epimorphism. Their composite is therefore a regular
epimorphism. It is equivariant because $B_T$ has trivial action.
Its top and bottom fibers have supports
\[
 P_T=G\times\mathsf{Th}_T,\qquad
 Q_T=G\times\mathsf{Rf}_T.
\]
These are disjoint c.e. subsets of $U$. Their partial identities
$e_{P_T},e_{Q_T}$ are restriction idempotents in $\mathcal R_G$.
Let $p_T^+=j_S(e_{P_T})$ and $p_T^-=j_S(e_{Q_T})$.

Decode unary-formula codes as in
Section~\ref{sec:integration-numerical}, and let
$\operatorname{sub}(a,n)=\beta^{-1}(\ulcorner\varphi_a(\nu_n)\urcorner)$
be the numerical substitution arrow. Define $s:U\times U\to S$ by
\[
 s((g,a),(h,n))=(g,\operatorname{sub}(a,n)).
\]
This map is total recursive and equivariant.

\begin{proposition}\label{prop:equivariant-logical-model}
The $\mathbf{Set}$-enriched category $\mathcal K$, with underlying
recursion category $\mathcal R_G$ and the chosen indexing arrow above, admits a
definability algebra with predicate objects $E_X$, total arrow $s$,
and distinguished restriction idempotents $e_{P_T},e_{Q_T}$.
Their images under $j_S$
give a definability algebra with predicate objects $P_X$.
Every arrow $JS\to\mathbf2$ in $\mathcal C$ is $\chi_D$ for a
complemented $D\in E(S)$, and there is no Boolean-algebra
homomorphism $B_T\to\mathbf2$ in $\mathcal C$.
\end{proposition}
\begin{proof}
A complemented restriction idempotent is $e_{G\times A}$ with $A$ recursive:
both $A$ and its complement are c.e., and dovetailing their
enumerations decides membership. For each $A$, choose the least
code $a$ of a unary formula strongly defining $A$, and put
$g_A(g,r)=(g,a)$.
This arrow is total recursive and equivariant. Substitution
gives
\[
 s(g_A\times1)((g,r),(h,n))
       =(g,\operatorname{sub}(a,n)).
\]
Strong numeralwise definability and consistency show that
its inverse images of $P_T$ and $Q_T$ are exactly the inverse
images under the second projection of $G\times A$ and
$G\times A^c$, respectively. These are the representation
equations~\eqref{eq:original-algebra-representation}.
Assigning $g_A$ to each pair
$(e_{G\times A},e_{G\times A^c})$ gives
a section of the representation projection in $\mathbf{Set}$.
Proposition~\ref{prop:algebra-comparison} then supplies the
section for the predicate objects $P_X$.

The object $\mathbf2=1+1$ has trivial action.
By the definition of arrows between equality PERs, every underlying
arrow $JU\to\mathbf2$ is tracked by a total recursive Boolean
function. Equivariance says its value at $(g,n)$ is independent
of $g$.  It therefore decides a recursive set of $n$, and its
two fibers are the images under $j$ of complementary restriction
idempotents.
Corollary~\ref{cor:integration-no-boolean-map} applies.
\end{proof}

The object $JU$ has no global elements: an equivariant arrow
$1\to JU$ would select an element fixed by every member of $G$,
whereas the action on $U$ is free.

The Boolean algebra $B_T$ and $\mathbf2$ both have trivial
action.  Hence their Boolean homomorphisms in $\mathcal C$ are
exactly their Boolean homomorphisms in
$\mathcal E_{\mathrm{rec}}$.  The usual numerical comparison
with complete decidable consistent extensions of $T$ therefore
remains valid. By Theorem~\ref{thm:integration-essential-equivalence},
the absence of these homomorphisms is equivalent to essential
incompleteness of $T$.

\section{Completions and introspective structures}
\label{sec:completion-introspection}

For a separable recursion category $\mathcal K_0$, adjoining a
restriction terminal object and splitting restriction idempotents
gives a restriction category whose total-map category $\mathcal D$
is finitely complete. Its exact completion provides effective
quotients. We next construct a free introspective extension of
a finitely complete category. Applied to the enriching base
$\mathscr V$, this extension changes the base of a definability
algebra and preserves its internal indexing and representation
operations. The unit has a finite-limit-preserving retraction.
The numerical examples
give categorical inseparability where every introspective functor
is naturally isomorphic to a constant functor with terminal value.
Arithmetic introspection supplies a representation operation after
realization in $\mathbf{Set}$; Section~\ref{sec:quotation-realization}
constructs this operation from quotation.

\subsection{Restriction terminal objects and exact completions}
\label{sec:completion-comparison}

We first construct finitely complete categories from the underlying
ordinary recursion category $\mathcal K_0$. There are two extensions
with a restriction terminal object. If $c:X\to X$ is a chosen total constant, so that
$c\pi_1=c\pi_2$, splitting $c$ gives the extension $\mathcal K_0^1$
of \cite[Sections~3--4]{RecursionCompletions}. Its objects are
$(A,1_A)$ and $(X,c)$, and its arrows satisfy
\[
 \mathcal K_0^1((A,e),(B,d))
   =\{f:A\to B\mid f=dfe\},
 \qquad \overline f^{\,1}=e\overline f.
\]
The object $(X,c)$ is restriction terminal. The inclusion of
$\mathcal K_0$ is fully faithful and preserves restrictions, zeros,
and binary restriction products. The original Turing morphism
also admits total indices with parameters in $\mathcal K_0^1$.
Functoriality here requires preservation of the chosen constant.

Without a chosen constant, a free restriction terminal extension is available
in the category of restriction categories and functors preserving
zeros and binary restriction products.

\begin{proposition}\label{prop:free-restriction-terminal}
Let $\mathcal K_0$ be a restriction category with restriction zeros
and binary restriction products. There is a fully faithful
restriction functor $J:\mathcal K_0\to\mathcal K_0^\top$ preserving
this structure, where $\mathcal K_0^\top$ has a restriction terminal
object $I$. Every structure-preserving functor
$F:\mathcal K_0\to\mathcal C$ into a category with a chosen
restriction terminal object extends uniquely to a
restriction-terminal-preserving functor on $\mathcal K_0^\top$
whose value at $I$ is the chosen restriction terminal object.
\end{proposition}
\begin{proof}
Adjoin one object $I$ and retain every old hom-set. For
$X\in\mathcal K_0$, define
\[
 \mathcal K_0^\top(X,I)=\{!e\mid e\in E(X)\},\qquad
 \mathcal K_0^\top(I,X)=\{0_{I,X}\},\qquad
 \mathcal K_0^\top(I,I)=\{0_I,1_I\},
\]
where $0_I\ne1_I$. Set $\overline{!e}=e$,
$\overline{0_{I,X}}=0_I$, and $\overline a=a$ for
$a\in\{0_I,1_I\}$. For $f:X\to Y$ in $\mathcal K_0$ and
$e\in E(Y)$, define
\[
 !e\,f= !\overline{ef}:X\longrightarrow I.
\]
The arrow $!0_X$ is the zero from $X$ to $I$; all zero arrows
absorb composition, and $1_I$ acts as identity. These rules cover
every composite involving $I$. Associativity for a composite
ending at $I$ follows from
$\overline{\overline{eg}f}=\overline{egf}$ for
$f:X\to Y$, $g:Y\to Z$, and $e\in E(Z)$; other new cases
are identity or zero composites. For $d,e\in E(X)$ the third
restriction equation reads $\overline{!d\,e}=de$.
The fourth reduces to $ef=f\overline{ef}$ for $f:X\to Y$
and $e\in E(Y)$. The remaining cases are the equations for
$0_I$ and $1_I$. Thus the old restriction structure extends.

Put $X\otimes I=I\otimes X=X$ and $I\otimes I=I$.
The projections from $X\otimes I$ are $1_X,!1_X$, and those
from $I\otimes I$ are both $1_I$. Retain the old products.
For arrows with a common source, the new pairings are
\[
 \langle f,!e\rangle=fe,\qquad
 \langle !e,f\rangle=fe,\qquad
 \langle !e,!d\rangle=!(ed).
\]
For source $I$, read $!e=e$ for $e\in\{0_I,1_I\}$;
an arrow $I\to X$ with $X\in\mathcal K_0$ is zero. Hence these
formulas also cover source $I$, including
$\langle1_I,1_I\rangle=1_I$. The projections give the
restriction-product equations; for example,
$\overline{fe}=\overline f\,e$ and
$!1_Xfe=!(\overline f\,e)$. Uniqueness follows from the
projection to the old factor, or from either identity projection
when both factors are $I$. For old factors and source $I$,
the unique pairing is zero.

The unique total arrow $X\to I$ is $!1_X$, and
$!1_Yf=!1_X\overline f$; the terminal arrow of $I$ is $1_I$.
If $1_{\mathcal C}$ is restriction terminal in $\mathcal C$,
the extension of $F$ is forced by
\[
 \widetilde F(I)=1_{\mathcal C},\qquad
 \widetilde F(!e)=!_{FX}F(e),\qquad
 \widetilde F(0_{I,X})=0_{1_{\mathcal C},FX}.
\]
These formulas preserve composition, restrictions, and pairings.
They also prove uniqueness.
\end{proof}

The extension $\mathcal K_0^\top$ has no total arrow from $I$ to
an old object. Thus $\mathcal K_0^1$ is the extension that preserves
the original Turing indexing for the new parameters, whereas
$\mathcal K_0^\top$ supplies a restriction terminal object for arbitrary
$\mathcal K_0$.

Write $\mathcal C$ for either restriction terminal extension and let
$K_r(\mathcal C)$ split its restriction idempotents. Its objects
are $(A,e)$ with $e=\overline e$, its arrows satisfy $f=dfe$,
and its identity at $(A,e)$ is $e$. Every original arrow
$f:X\to Y$ is retained and determines the span of total maps
\[
\begin{tikzcd}
 X & (X,\overline f)
   \arrow[l, hook, "\overline f"'] \arrow[r, "f"] & Y
\end{tikzcd}
\]
in $\mathcal D=\operatorname{Tot}K_r(\mathcal C)$.
Thus passing to total maps records each partial arrow $f$ together
with the splitting of its restriction idempotent $\bar f$.

For the original enrichment $\mathcal K$, splittings of idempotents
admit the following description in $\mathscr V$.
For global idempotents $e:A\to A$ and $d:B\to B$ of
$\mathcal K_0$, define a hom-object in $\mathscr V$ by
\[
 H_{e,d}=\operatorname{Eq}\bigl(1_{H(A,B)},
             f\longmapsto dfe\bigr).
\]
Here multiplication by $e$ and $d$ uses their representing points
in the hom-objects. For a third global idempotent $k:C\to C$,
enriched composition restricts to
$H_{d,k}\times H_{e,d}\to H_{e,k}$, and the identity at $(A,e)$
is the point represented by $e$. Preservation of equalizers by
$\Gamma$ identifies its ordinary hom-set with
\[
 \Gamma H_{e,d}\cong
 \{f\in\mathcal K_0(A,B)\mid f=dfe\}.
\]
For total idempotents, restriction is induced by
$f\mapsto e\bar f$; for restriction idempotents this formula
reduces to $f\mapsto\bar f$. Indeed, $fe=f$ implies
$\bar f e=e\bar f$ in the first case and $\bar f\leq e$ in
the second. The restriction equations therefore restrict to the
equalizer hom-objects. Binary restriction products are
$(A,e)\otimes(B,d)=(A\otimes B,e\otimes d)$, with projections
$\pi_1(e\otimes d)$ and $\pi_2(e\otimes d)$ and pairing
induced by the original pairing operation. Thus these constructions
enrich the total-idempotent and restriction-idempotent splittings
over $\mathscr V$. The chosen-constant extension is the full
enriched subcategory on $(A,1_A)$ and $(X,c)$, with products
transported along the isomorphisms of
\cite[Section~3]{RecursionCompletions}.
Finite-limit-preserving change of base preserves these hom-object
equalizers for the chosen split objects.

An object $A$ is \emph{separable} if its diagonal has a
restriction inverse $\delta_A:A\otimes A\to A$ satisfying
\[
 \delta_A\Delta_A=1_A,\qquad
 \Delta_A\delta_A=\overline{\delta_A}.
\]
Assume every object of $\mathcal K_0$ is separable. This property
passes to $\mathcal K_0^1$ and to restriction-idempotent splittings
\cite[Section~4]{RecursionCompletions}. It also passes to
$\mathcal K_0^\top$, since the old products are retained and
$\Delta_I=1_I$ has restriction inverse $1_I$.
The category $\mathcal D$
has finite limits: it has a terminal object and binary products,
and the equalizer of total $f,g:A\to B$ is the splitting of
$e_{f,g}=\overline{\delta_B\langle f,g\rangle}$.
Indeed, $fe_{f,g}=ge_{f,g}$. If a total $h:Z\to A$ satisfies
$fh=gh$, then
$\delta_B\langle f,g\rangle h=fh$ is total, so
$e_{f,g}h=h$; this gives the unique factorization through the
splitting. A restriction functor preserving binary products
preserves the restriction inverses of diagonals and these
equalizers. If it also preserves the restriction terminal object,
its induced functor on $\mathcal D$ preserves all finite limits.
The functors induced by the two terminal extensions above have
these preservation properties.

For small $\mathcal D$, the free exact completion of
Carboni and Celia Magno \cite{CarboniCeliaMagno1982} gives
\[
 \mathcal D\longrightarrow\operatorname{Ex}_{\mathrm{lex}}
 (\mathcal D),\qquad
 \operatorname{Reg}(\operatorname{Ex}_{\mathrm{lex}}(\mathcal D),
                    \mathcal E)
 \simeq \operatorname{Lex}(\mathcal D,\mathcal E)
\]
for every exact category $\mathcal E$. Here $\operatorname{Lex}$
and $\operatorname{Reg}$ denote the categories of finite-limit-
preserving and regular functors, respectively. Equivalently, one
may first take the free regular completion and then its exact
completion \cite{CarboniVitale}:
\[
 \operatorname{Ex}_{\mathrm{lex}}(\mathcal D)
 \simeq
 \operatorname{Ex}_{\mathrm{reg}}
       (\operatorname{Reg}_{\mathrm{lex}}(\mathcal D)).
\]
When $\mathcal D$ is already regular, the construction in
\cite[Section~6]{RecursionCompletions} instead uses
$\operatorname{Ex}_{\mathrm{reg}}(\mathcal D)$, whose universal
property is
\[
 \operatorname{Reg}(\operatorname{Ex}_{\mathrm{reg}}(\mathcal D),
                    \mathcal E)
 \simeq \operatorname{Reg}(\mathcal D,\mathcal E).
\]
Its functoriality requires the induced functor on total maps to
preserve regular epimorphisms. In the numerical category
$\mathcal D_{\mathrm{rec}}$, every regular epimorphism splits,
so every finite-limit-preserving functor into a regular category
has this property. The $\mathrm{lex}$ construction supplies an
exact completion when regularity has not been established.

\subsection{Introspective structures and free extensions}
\label{sec:introspective-extension}

For the free-extension and change-of-base constructions, bases carry
chosen terminal objects and pullbacks, and base functors preserve
these choices strictly. Products and
equalizers are formed from these chosen limits. Work in a universe
in which the bases and enriched object sets under consideration
are small, and form categories of these structures in a larger
universe. In the numerical examples, $\mathbf{Set}$ denotes sets
in a smaller universe.

Let $\mathbf{Lex}_{\mathrm{s}}$ denote the category of these
finite-limit categories and functors. For a finite-limit category
$A$, an internal finite-limit category $C$ has objects of objects
and arrows in $A$, internal category operations, and chosen internal
finite limits. Write $\operatorname{Glob}_A(C)$ for its category of
global objects and arrows. For global objects $a,c:1\to C_0$,
the hom-object $C(a,c)$ is the pullback of
$(\operatorname{src},\operatorname{tgt}):C_1\to C_0\times C_0$
along $(a,c):1\to C_0\times C_0$.
In particular, $C(1_C,c)$ belongs to $A$ for each global
object $c$ of $C$.

Following Ramesh, a \emph{strict introspective theory} consists of
$A\in\mathbf{Lex}_{\mathrm{s}}$, an internal finite-limit category
$C$ in $A$, a strict finite-limit-preserving functor
\[
 \mathsf S:A\longrightarrow\operatorname{Glob}_A(C),
\]
and a natural transformation
\[
 N:\mathrm{id}_A\Longrightarrow\Box_A,
 \qquad \Box_A(X)=C(1_C,\mathsf SX).
\]
Thus $N_X:X\to C(1_C,\mathsf SX)$ is an arrow of $A$.
For an arrow $f:X\to Y$ of $A$, the morphism
\[
 \Box_A(f)=C(1_C,\mathsf Sf):
 C(1_C,\mathsf SX)\longrightarrow C(1_C,\mathsf SY)
\]
is induced by postcomposition with $\mathsf Sf$.
Naturality of $N$ is the equation $\Box_A(f)N_X=N_Yf$.
A homomorphism $(A,C,\mathsf S,N)\to(A',C',\mathsf S',N')$ is a strict
finite-limit-preserving functor $H:A\to A'$ satisfying
\[
 H[C]=C',\qquad
 H_C \mathsf S=\mathsf S'H,\qquad H(N_X)=N'_{HX},
\]
where $H_C:\operatorname{Glob}_A(C)\to
\operatorname{Glob}_{A'}(C')$ applies $H$ to global objects and
arrows. Preservation of the chosen limits identifies the codomains
in the last equation. These are Ramesh's Definitions~5.8 and~5.10
\cite[pp.~93--94]{Ramesh}. Denote this category by
$\mathbf{Int}_{\mathrm{s}}$ and its underlying-category functor by
$U:\mathbf{Int}_{\mathrm{s}}\to\mathbf{Lex}_{\mathrm{s}}$.

\begin{proposition}\label{prop:free-introspective-extension}
The functor $U$ has a left adjoint $L$. For every $A$, its unit
$\eta_A:A\to UL(A)$ has a strict finite-limit-preserving retraction.
Consequently $\eta_A$ is faithful and reflects isomorphisms.
\end{proposition}
\begin{proof}
Strict introspective theories are models of an essentially algebraic
theory extending the theory of strict finite-limit categories
\cite[p.~94]{Ramesh}. Add to its signature a name for every object
and arrow of $A$, with equations for their domains, codomains,
identities, compositions, and chosen finite limits. The initial model
of this presentation is $L(A)$; the named objects and arrows define
$\eta_A$. Interpreting these names in an introspective theory $I$
is precisely giving a strict finite-limit-preserving functor
$A\to UI$. Hence
\[
 \mathbf{Int}_{\mathrm{s}}(L(A),I)
 \cong\mathbf{Lex}_{\mathrm{s}}(A,UI).
\]
The unique extension of $\eta_Bf$ defines $L(f)$ for
$f:A\to B$. Uniqueness proves the functor laws and naturality of
the displayed bijection.

Let $R(A)$ have underlying category $A$ and the terminal internal
category: both its object of objects and its object of arrows are
$1_A$, and all operations are the unique possible arrows. Take $\mathsf S$
to be the constant functor at its terminal object and $N_X$ to be
the unique arrow from $X$ to the terminal object $\Box_A X$.
These choices give an introspective theory and are preserved by
every strict finite-limit-preserving functor. Thus $R$ is a functor
with $UR=\mathrm{id}$.

The identity $A\to UR(A)$ extends uniquely to a homomorphism
$\epsilon_A:L(A)\to R(A)$. Put $r_A=U\epsilon_A$. Then
$r_A\eta_A=\mathrm{id}_A$. Applying $r_A$ shows that $\eta_A$
reflects equality of arrows and invertibility.
\end{proof}

\subsection{Change of base for definability algebras}
\label{sec:definability-base-change}

The completions in Section~\ref{sec:completion-comparison} act on
the underlying recursion category $\mathcal K_0$. Change of base
acts on the ambient category $\mathscr V$ containing the hom-objects,
predicate objects, and representation operation.

For the chosen presentations in this section, take the objects
$E_X$, $\mathsf T(X,Y)$, $C_U$, and $R$ from the fixed finite-limit
constructions above. If $F:\mathscr V\to\mathscr W$ is a strict
finite-limit-preserving functor and $D$ is a definability algebra
over $\mathscr V$, apply $F$ to its hom-objects, predicate objects,
and structure arrows. Denote these data by $F_*D$ and its
$\mathscr W$-enriched category by $F_*\mathcal K$. Its objects
and chosen products are those of $\mathcal K$, and its hom-objects
are $FH(X,Y)$. The underlying ordinary category has hom-sets
\[
 (F_*\mathcal K)_0(X,Y)=\mathscr W(1,FH(X,Y)).
\]
Applying $F$ to global arrows defines an ordinary functor
\[
 \iota_F:\mathcal K_0\longrightarrow(F_*\mathcal K)_0,
 \qquad \widehat f\longmapsto F(\widehat f).
\]

\begin{lemma}\label{lem:definability-base-change}
The structure $F_*D$ is a definability algebra over $\mathscr W$.
The functor $\iota_F$ preserves restriction, the chosen restriction
products, zeros, and joins of restriction idempotents. If $F$ is faithful, so is
$\iota_F$, and $1_U\ne0_{U,U}$ implies the same inequality in
$(F_*\mathcal K)_0$. If $F$ is fully faithful, $\iota_F$ is an
isomorphism of ordinary categories.
\end{lemma}
\begin{proof}
The enriched category, restriction, product, and lattice axioms
are equations of structure arrows, so application of $F$ preserves
them. Zero absorption by the restriction product also follows
equationally: for $h=\langle f,0\rangle$,
$\bar h=\bar f\,\bar0=0$, whence $h=h\bar h=0$.
Thus the prodominical equations hold. The original isomorphisms
between objects remain isomorphisms, so $(F_*\mathcal K)_0$ is an isotype.

Preservation of the defining limits gives
\[
 E_X^{F}=F(E_X),\quad
 \mathsf T^{F}(X,Y)=F(\mathsf T(X,Y)),\quad
 C_U^{F}=F(C_U),\quad R^{F}=F(R).
\]
The morphism $F(j_X)$ is monic, and the predicate actions and their
compatibility with $j$ satisfy the required equations. The equation
$F(q_R)F(\sigma)=1_{F(C_U)}$ makes $F(\sigma)$ the representation
section. Its representation operation is
$\rho^F=F(\rho):F(C_U)\longrightarrow F\mathsf T(U,U)$.

The indexing equation becomes
\[
 F(\operatorname{ev}_{\tau})F(\Lambda)
     =1_{FH(U\otimes U,U)}.
\]
For every global arrow $a:1\to FH(U\otimes U,U)$, the composite
$F(\Lambda)a:1\to F\mathsf T(U,U)$ therefore gives a total index
for $a$. Write $\tau_F:U\otimes U\to U$ for the ordinary arrow
represented by $F(\tau)$. For an arbitrary parameter object $X$,
choose an isomorphism $i:X\to U$. Given $f:X\otimes U\to U$,
apply this indexing construction to $f(i^{-1}\otimes1_U)$ to
obtain a total $g:U\to U$. Then $gi:X\to U$ is total and
$\tau_F((gi)\otimes1_U)=f$. Thus $\tau_F$ is a Turing morphism
in $(F_*\mathcal K)_0$.
Likewise, every global pair of complementary restriction idempotents in
$(F_*\mathcal K)_0$ factors through $F(C_U)$, so $F(\sigma)$ supplies
its representative. These verifications apply to all global
arrows of the new hom-objects.

The preservation statements for $\iota_F$ follow from its
definition. If $F$ is faithful, equality of the images of two
global arrows implies their equality, proving faithfulness of
$\iota_F$ and preservation of $1_U\ne0_{U,U}$. If $F$ is fully
faithful, its maps
\[
 \mathscr V(1,H(X,Y))\longrightarrow
 \mathscr W(F1,FH(X,Y))=\mathscr W(1,FH(X,Y))
\]
are bijections. Since $\iota_F$ is the identity on objects, it is
then an isomorphism of ordinary categories.
\end{proof}

In particular, for small $\mathscr V$ the canonical embedding
$\mathscr V\to\operatorname{Ex}_{\mathrm{lex}}(\mathscr V)$ is
fully faithful and preserves finite limits
\cite[Section~2]{CarboniCeliaMagno1982}. With compatible choices
of finite limits, change of base along this embedding gives a
definability algebra over an exact category with the same underlying
ordinary recursion category up to the isomorphism above.

Change of base also applies to morphisms of definability algebras.
With the chosen strict limits,
$(GF)_*=G_*F_*$ and $(1_{\mathscr V})_*=1$.

\subsection{The relative adjunction}
\label{sec:relative-introspective-adjunction}

Define $\mathbf{Def}$ to have objects $(\mathscr V,D)$, where
$\mathscr V\in\mathbf{Lex}_{\mathrm{s}}$ and $D$ is a definability
algebra over $\mathscr V$. A morphism
$(\mathscr V,D)\to(\mathscr W,D')$ is a pair $(F,\alpha)$,
where $F:\mathscr V\to\mathscr W$ is strict and preserves finite
limits, and $\alpha:F_*D\to D'$ is a morphism of definability
algebras over $\mathscr W$ as defined in Section~\ref{sec:definability-morphisms}.
Composition and identity are
\[
 (G,\beta)(F,\alpha)=(GF,\beta\,G_*(\alpha)),
 \qquad 1_{(\mathscr V,D)}=(1_{\mathscr V},1_D).
\]
An object of $\mathbf{IntDef}$ is $(I,D)$, where $I$ is a strict
introspective theory and $D$ is a definability algebra over $UI$.
Its morphisms are the pairs above whose base functors preserve
the introspective structures.
Write $U_D:\mathbf{IntDef}\to\mathbf{Def}$ for the functor
forgetting the introspective structure.

\begin{theorem}\label{thm:relative-introspective-adjunction}
There is an adjunction
\[
 L_D:\mathbf{Def}\rightleftarrows\mathbf{IntDef}:U_D,
 \qquad L_D\dashv U_D,
\]
with
\[
 L_D(\mathscr V,D)=
 \bigl(L(\mathscr V),(\eta_{\mathscr V})_*D\bigr).
\]
Its unit has a retraction in $\mathbf{Def}$. On underlying
recursion categories it induces a faithful functor with a
structure-preserving retraction. If $1_U\ne0_{U,U}$ in $D$,
then the distinguished predicates of $(\eta_{\mathscr V})_*D$
have no separator given by a complemented restriction idempotent.
\end{theorem}
\begin{proof}
A morphism from $L_D(\mathscr V,D)$ to $(I,D')$ consists of an
introspective homomorphism $H:L(\mathscr V)\to I$ and an algebra
morphism
\[
 (UH)_*\bigl((\eta_{\mathscr V})_*D\bigr)\longrightarrow D'.
\]
The adjunction $L\dashv U$ identifies $H$ uniquely with
$F=UH\eta_{\mathscr V}:\mathscr V\to UI$. Since
$(UH)_*(\eta_{\mathscr V})_*D=F_*D$, the algebra morphism is
exactly the second component of a morphism
$(\mathscr V,D)\to U_D(I,D')$. This gives the natural hom-set
bijection. On a morphism $(F,\alpha)$, the left adjoint is
$(L(F),(\eta_{\mathscr W})_*\alpha)$, using naturality of $\eta$.

The unit is $(\eta_{\mathscr V},1_{(\eta_{\mathscr V})_*D})$.
The retraction $r_{\mathscr V}$ of
Proposition~\ref{prop:free-introspective-extension} satisfies
$(r_{\mathscr V})_*(\eta_{\mathscr V})_*D=D$. Paired with $1_D$,
it therefore retracts this unit. On underlying recursion
categories the induced functors are
\[
 \mathcal K_0
 \xrightarrow{\ \iota_{\eta_{\mathscr V}}\ }
 ((\eta_{\mathscr V})_*\mathcal K)_0
 \xrightarrow{\ \iota_{r_{\mathscr V}}\ }\mathcal K_0.
\]
Their composite is the identity because
$r_{\mathscr V}\eta_{\mathscr V}=1_{\mathscr V}$.
Lemma~\ref{lem:definability-base-change} gives faithfulness and
preservation of nontriviality. The final assertion follows from
Corollary~\ref{cor:definability-algebra-inseparability}.
\end{proof}

\subsection{Definability without nonconstant introspection}
\label{sec:introspection-separation}

The numerical definability algebras give a separation from
nonconstant introspection on the same base category. We first prove a
restriction on finite-limit-preserving functors with the
L\"ob property.

Let $\mathcal A$ have finite limits and terminal object $1$.
A finite-limit-preserving endofunctor $\Box:\mathcal A\to\mathcal A$
has the \emph{L\"ob property} if, for every
$f:\Box X\to X$, there is an arrow $x:1\to X$ satisfying
\[
 f\Box(x)u=x,
 \qquad u:1\xrightarrow{\cong}\Box1,
\]
where $u$ is the unique isomorphism between the indicated terminal
objects \cite[Definition~4.1]{Ramesh}.
An initial object $0$ is \emph{strict} if every arrow into $0$
is an isomorphism. An object is \emph{subterminal} if its unique
arrow to $1$ is monic.

\begin{theorem}\label{thm:two-valued-lob-collapse}
Let $\mathcal A$ have finite limits and a strict initial object
$0$, with $0\not\cong1$. Suppose that its only subobjects of
$1$ are $0$ and $1$. Every finite-limit-preserving endofunctor
$\Box:\mathcal A\to\mathcal A$ with the L\"ob property is
naturally isomorphic to the constant functor with value $1$.
\end{theorem}
\begin{proof}
Since $\Box$ preserves monomorphisms and the terminal object,
$\Box0$ is subterminal. Thus $\Box0\cong0$ or $\Box0\cong1$.
In the first case an isomorphism $\Box0\to0$ and the L\"ob
property give an arrow $1\to0$. Strictness would make this
arrow an isomorphism, contrary to $0\not\cong1$.
Hence $\Box0\cong1$.

For each $X$, strictness makes the projection $0\times X\to0$
an isomorphism. Applying $\Box$ and using its product comparison
shows that the projection
$\Box0\times\Box X\longrightarrow\Box0$
is an isomorphism. Since $\Box0$ is terminal, the unique arrow
$\Box X\to1$ is therefore an isomorphism. These unique arrows
form the required natural isomorphism.
\end{proof}

For an introspective structure, the endofunctor $\Box$ defined
in Section~\ref{sec:introspective-extension} preserves finite limits
\cite[Definition~3.2]{Ramesh} and has the L\"ob property
\cite[Theorem~4.19 and Corollary~4.21]{Ramesh}.

\begin{corollary}\label{cor:two-valued-introspection-collapse}
Under the hypotheses on $\mathcal A$ in
Theorem~\ref{thm:two-valued-lob-collapse}, every introspective
structure $(\mathcal A,\mathcal C,\mathsf S,N)$ has
$\mathsf S$ naturally isomorphic to the constant functor
with value $1_{\mathcal C}$, and $\Box$ naturally isomorphic
to the constant functor with value $1$.
\end{corollary}
\begin{proof}
The theorem gives $\Box0\cong1$. A global point of this
internal hom is an arrow $1_{\mathcal C}\to \mathsf S0$ in
$\operatorname{Glob}(\mathcal C)$. Since $\mathsf S$ preserves finite
limits, $\mathsf S0$ is subterminal. Its point therefore makes it
terminal. Applying $\mathsf S$ to the projection $0\times X\to0$
now shows that $\mathsf S0\times \mathsf SX\to \mathsf S0$ is an isomorphism, so
$\mathsf SX$ is terminal for every $X$. The unique arrows from these
objects to $1_{\mathcal C}$ give the natural isomorphism.
\end{proof}

The internal category $\mathcal C$ may have nonterminal objects;
the corollary says that the image of $\mathsf S$ consists of terminal
objects. These are the constant-functor introspective structures
of \cite[Example~2.4]{Ramesh}.

\begin{corollary}\label{cor:numerical-introspection-separation}
Let $T$ be a consistent c.e. extension of PA. Its numerical
definability algebra over $\mathbf{Set}$ has restriction idempotents as predicates,
and its Lindenbaum presentation belongs to
$\mathcal E_{\mathrm{rec}}$.
Every introspective structure on any of
$\mathbf{Set}$, $\mathcal E_{\mathrm{rec}}$, and $\mathcal P$,
with internal category $\mathcal C$, has $\mathsf S$ naturally
isomorphic to the constant functor with value
$1_{\mathcal C}$.
\end{corollary}
\begin{proof}
The definability algebra is supplied by
Corollaries~\ref{cor:numerical-definability-algebras}
and~\ref{cor:arithmetic-representation-incompleteness}.
Since $T$ is c.e., its Lindenbaum PER and presentation belong
to $\mathcal E_{\mathrm{rec}}$, as established in
Section~\ref{sec:integration-numerical}.

The category $\mathbf{Set}$ satisfies the hypotheses of
Theorem~\ref{thm:two-valued-lob-collapse}. In each of the two
PER categories, the empty PER is strict initial: an arrow into
it forces the source support to be empty. The singleton equality
PER is terminal and is not initial. If a monomorphism
$R\rightarrowtail1$ has nonempty support, choose a supported
representative $n$. The constant recursive function with value
$n$ gives a section $1\to R$. The monomorphism is consequently
an isomorphism. Thus the only subobjects of $1$ are $0$ and $1$.
Corollary~\ref{cor:two-valued-introspection-collapse} applies.
\end{proof}

Here $\mathbf{Set}$ is the ambient category containing the
hom-objects and the representation operation. The category
$\mathcal E_{\mathrm{rec}}$ is the exact completion providing
the numerical Boolean presentation. The corollary excludes
nonconstant introspection on both categories.

There is also an obstruction to identifying the embedding of restriction idempotents
with the natural transformation $N$.

\begin{proposition}\label{prop:domain-embedding-not-introspection}
Let $\mathcal A$ satisfy the hypotheses of
Theorem~\ref{thm:two-valued-lob-collapse} and carry an
introspective structure. If $E$ has distinct global points and
$j:E\rightarrowtail P$ is monic, then $N_E:E\to\Box E$
is not monic, and no arrow $a:\Box E\to P$ satisfies
$aN_E=j$.
\end{proposition}
\begin{proof}
Choose distinct $b,t:1\to E$. The object $\Box E$ is terminal,
so $N_Eb=N_Et$. This proves that $N_E$ is not monic.
If $aN_E=j$, then $jb=jt$, contradicting monicity of $j$.
\end{proof}

For the numerical definability algebra, take $E=E_U$ and
$j=j_U$. Its lower and upper bounds are distinct global
points. Proposition~\ref{prop:domain-embedding-not-introspection}
therefore prevents the identification of $N_{E_U}$ with the
embedding of restriction idempotents, even after an isomorphism of their codomains.

\subsection{Numerical realization of arithmetic introspection}
\label{sec:quotation-realization}

A finite-limit-preserving interpretation of the arithmetic
introspective structure sends its internal category and quotation
maps to $\mathbf{Set}$. From these interpretations we construct
the representation operation for the numerical $\mathbf{Set}$-enrichment.
The construction uses disjoint subobjects describing
the two possible outputs of a program. For a total Boolean
program, their inverse images under its input map are complementary.

\begin{lemma}\label{lem:realized-disjoint-subobjects}
Let $(\mathcal B,\mathcal C,\mathsf S,N)$ be an introspective
structure, and suppose that $\mathcal B$ has an initial object
$0_{\mathcal B}$. Let $F:\mathcal B\to\mathcal E$ preserve finite
limits, where $\mathcal E$ is finitely complete and has a strict initial object
$0_{\mathcal E}$ and distributive subobject lattices. Suppose that
$F(\Box0_{\mathcal B})\cong0_{\mathcal E}$.
For disjoint monomorphisms $m_\pm:M_\pm\rightarrowtail Z$ in
$\mathcal B$, put
$Q_\pm=(FN_Z)^*[F(\Box m_\pm)]$ in
$\operatorname{Sub}_{\mathcal E}(FZ)$.
Then $[Fm_\pm]\leq Q_\pm$ and $Q_+\wedge Q_-=0$.
For every $h:Y\to FZ$ satisfying
$h^*[Fm_+]\vee h^*[Fm_-]=1_Y$, one has
$h^*Q_\pm=h^*[Fm_\pm]$.
\end{lemma}
\begin{proof}
Naturality gives the commutative square
\[
\begin{tikzcd}
 FM_\pm \arrow[r,"F(N_{M_\pm})"] \arrow[d,hook,"Fm_\pm"']
   & F(\Box M_\pm) \arrow[d,hook,"F(\Box m_\pm)"]\\
 FZ \arrow[r,"FN_Z"'] & F(\Box Z).
\end{tikzcd}
\]
Its factorization through the pullback defining $Q_\pm$ gives
$[Fm_\pm]\leq Q_\pm$. Since $F$ and $\Box$ preserve finite
limits, the intersection of $F(\Box m_+)$ and $F(\Box m_-)$ is
$F(\Box0_{\mathcal B})$. The hypothesis on this object gives
$Q_+\wedge Q_-=0$.

Write $D_\pm=h^*[Fm_\pm]$ and $A_\pm=h^*Q_\pm$.
Then $D_\pm\leq A_\pm$, $A_+\wedge A_-=0$, and
$D_+\vee D_-=1_Y$. Distributivity gives
$A_+=A_+\wedge(D_+\vee D_-)=D_+$; the other equality follows
in the same way.
\end{proof}

We use the arithmetical form of Ramesh's syntactic constructions
\cite[Constructions~2.9--2.10, pp.~49--53]{Ramesh}.
Let $\mathcal B$ be the exact completion of the regular
$\Sigma_1$ syntactic category of PA. It has the following PER
presentation. An object is a $\Sigma_1$ formula $R(x,x')$ whose
symmetry and transitivity PA proves. An arrow $R\to R'$ is a
$\Sigma_1$ formula $H(x,y)$ for which PA proves the universal
closures of
\begin{align*}
 H(x,y)&\ \Longrightarrow\ R(x,x)\wedge R'(y,y),\\
 R(x,x')\wedge R'(y,y')\wedge H(x,y)
   &\ \Longrightarrow\ H(x',y'),\\
 R(x,x)&\ \Longrightarrow\ \exists y\,H(x,y),\\
 H(x,y)\wedge H(x,y')&\ \Longrightarrow\ R'(y,y').
\end{align*}
Arrows are identified when PA proves their defining formulas
equivalent. Composition is relational composition. Write
$\mathcal N$ for the equality PER defined by $x=y$,
$0_{\mathcal B}$ for the empty PER, and $1_{\mathcal B}$ for
equality on the singleton defined by $x=0$.

Let $T=\mathrm{PA}+\Gamma$, where $\Gamma$ is a c.e. collection
of axioms in an effective classical language extending that of
arithmetic. A proof includes finite computation records verifying
its syntactic checks and its uses of axioms from $\Gamma$.
This presentation gives a
primitive recursive proof-verification relation and a
PA-verifiable translation of PA proofs into $T$-proofs.
Define the internal category $\mathcal C_T$ in $\mathcal B$
by the same syntactic construction, allowing arbitrary formulas
for the PERs and functional relations and using $T$-provability
in the defining conditions. Its object of objects is the PER of
valid PER-formula codes, with two codes related when $T$ proves
their formulas equivalent. Its object of arrows is the PER of
triples $(R,H,R')$ satisfying the functional-relation conditions;
two triples are related when $T$ proves their corresponding
component formulas equivalent.
The defining relations on these codes are $\Sigma_1$, and the
exact completion supplies the quotients. Identities are represented
by the PER formulas, and composition is relational composition.
For arrows $H:R\to V$ and $G:R'\to V$, their pullback is represented by
\[
 P((x,y),(x',y'))\ \Longleftrightarrow\
 R(x,x')\wedge R'(y,y')\wedge
 \exists z\bigl(H(x,z)\wedge G(y,z)\bigr).
\]
Saturation makes equality of the two images equivalent to the
existence of the common representative $z$. The first projection
has graph $P((x,y),(x,y))\wedge R(x,r)$, and the second has
graph $P((x,y),(x,y))\wedge R'(y,r')$.
For compatible arrows $L:W\to R$ and $K:W\to R'$, the graph
$L(w,x)\wedge K(w,y)$ defines the mediating arrow.
Compatibility places its values in the support of $P$;
the two projection equations give uniqueness modulo $P$.
Replacing any component formula by a $T$-provably equivalent
formula gives a $T$-provably equivalent construction. These
operations therefore descend to the PER quotients and give the
internal finite-limit structure, with terminal PER equality on
$\{0\}$. The constructions on formulas and proofs are primitive
recursive, and PA verifies their correctness. They therefore
define structure arrows in $\mathcal B$, as in Ramesh's
Construction~2.10.

Regard each base formula and arrow as the same formula modulo
$T$-provability. This gives the finite-limit-preserving functor
$\mathsf S_T:\mathcal B\to
\operatorname{Glob}_{\mathcal B}(\mathcal C_T)$.
Thus $\mathsf S_TX$ is a global object of $\mathcal C_T$, and
$\Box_TX=\mathcal C_T(1_{\mathcal C_T},\mathsf S_TX)$ is an
object of $\mathcal B$. For a base PER $R$, a supported
representative $n$ determines the constant numeral arrow into
$\mathsf S_TR$, represented by
$(z=0)\wedge R(\overline n,y)$, where $z$ is the terminal input.
Write $\operatorname{Pr}_T$ for the arithmetic proof predicate
from the chosen proof coding, and let
$\ulcorner\alpha(\dot{\vec x})\urcorner$ denote the code obtained
by substituting the numerals for $\vec x$ into a fixed formula $\alpha$.
Formalized $\Sigma_1$-completeness and the proof translation into
$T$ give, for each $\Sigma_1$ formula $\alpha$, the PA theorem
$\alpha(\vec x)\Rightarrow
\operatorname{Pr}_T(\ulcorner\alpha(\dot{\vec x})\urcorner)$.
Applied to $R$, this shows that the assignment respects $R$ and defines
$N^T_R:R\to\Box_TR$ in $\mathcal B$.
Applied to the functional relation $H(x,y)$ of an arrow, the
same PA theorem proves that applying $\mathsf S_TH$ to the
numeral arrow for $x$ gives the numeral arrow for $y$ whenever
$H(x,y)$ holds. These PA-provable equations give naturality of
$N^T:\mathrm{id}_{\mathcal B}\Rightarrow\Box_T$.
These constructions give an introspective structure
$(\mathcal B,\mathcal C_T,\mathsf S_T,N^T)$.
Chosen finite-limit presentations can be made strict by
\cite[Theorem~5.9]{Ramesh}.

Interpret base formulas in the standard model of PA. This gives
a finite-limit-preserving functor $F:\mathcal B\to\mathbf{Set}$:
it sends $R$ to the quotient of
$\{n\in\mathbb N:\mathbb N\models R(n,n)\}$ by $R$, and sends a functional
relation to its induced function on equivalence classes.
Interpretation preserves the finite-limit constructions on
these quotients. In particular, $F\mathcal N=\mathbb N$ and
$F0_{\mathcal B}=\varnothing$.
Applying $F$ to the internal-category diagram gives the ordinary
category $F(\mathcal C_T)$ of $T$-provable PERs and functional
relations, modulo the equivalences just defined.
The object $\mathsf S_T0_{\mathcal B}$ is the empty PER in the
internal syntactic category. A point of $F(\Box_T0_{\mathcal B})$
is a class of $T$-provably functional relations from its terminal
object to that empty PER. Such a relation exists precisely when
$T\vdash\bot$. Thus consistency of $T$ gives
$F(\Box_T0_{\mathcal B})=\varnothing$.

\begin{theorem}\label{thm:introspection-numerical-definability}
Let $T=\mathrm{PA}+\Gamma$ be consistent, where $\Gamma$ is c.e.
in an effective classical language extending arithmetic.
The introspective structure
$(\mathcal B,\mathcal C_T,\mathsf S_T,N^T)$ and its standard
realization $F$ determine a representation section for the
numerical recursion category $\mathcal R_{\mathrm{num}}$ with
$U=S=\mathbb N$, predicate objects $P_X=E_X$, and the numeral
substitution of Section~\ref{sec:integration-numerical}. Its distinguished
restriction idempotents are $a_T^+=e_{\mathsf{Th}_T}$ and $a_T^-=e_{\mathsf{Rf}_T}$,
where $e_A$ denotes the partial identity on $A$.
Consequently these data form a definability algebra over
$\mathbf{Set}$.
\end{theorem}
\begin{proof}
Fix the acceptable program numbering used for the Turing
morphism $\tau:U\otimes U\rightharpoonup U$ in
$\mathcal R_{\mathrm{num}}$.
Let $M_b(p,x)$, for $b\in\{0,1\}$, be a $\Sigma_1$ formula
asserting that program $p$ on input $x$ terminates with output
$b$. Its computation relation is deterministic, and PA proves
that $M_0(p,x)$ and $M_1(p,x)$ are disjoint.
These formulas therefore define disjoint subobjects
$m_b:\mathcal M_b\rightarrowtail\mathcal N\times\mathcal N$
in $\mathcal B$.

Put $Z=\mathcal N\times\mathcal N$ and define $Q_b$ by the
pullback
\[
\begin{tikzcd}
 Q_b \arrow[r] \arrow[d,hook]
   & F(\Box_T\mathcal M_b) \arrow[d,hook,"F(\Box_Tm_b)"]\\
 \mathbb N^2 \arrow[r,"FN^T_Z"'] & F(\Box_TZ).
\end{tikzcd}
\]
For a standard pair $(p,n)$, membership in $Q_b$ says that its
numeral arrow factors through the interpretation of $\mathsf S_Tm_b$
in $F(\mathcal C_T)$. A proof of $M_b(\overline p,\overline n)$
gives such a factorization by the constant numeral arrow into
the interpretation of $\mathsf S_T\mathcal M_b$.
Conversely, write $G(z)$ for the graph of a factorization after
substitution of the unique terminal input. Its defining proofs give
$T\vdash\exists z\,G(z)$ and
$T\vdash\forall z\,(G(z)\Rightarrow
(M_b(z)\wedge z=(\overline p,\overline n)))$, with pairs
encoded by the fixed tuple coding. Hence
$T\vdash M_b(\overline p,\overline n)$.
Thus $(p,n)\in Q_b$ if and only if
$T\vdash M_b(\overline p,\overline n)$.

Let $p$ be an index of an actual total Boolean function. The map
$h_p:\mathbb N\to\mathbb N^2$, given by $h_p(n)=(p,n)$, has
complementary inverse images of $Fm_0$ and $Fm_1$.
Lemma~\ref{lem:realized-disjoint-subobjects} gives
$h_p^*Q_b=h_p^*[Fm_b]$.
The coverage used here holds in $\mathbf{Set}$ because the
program terminates with output $0$ or $1$ on every input.
It therefore applies to every total Boolean program, including
those whose totality $T$ does not prove.

The fixed recursive tuple coding identifies the underlying set
of $U\otimes U$ with $\mathbb N^2=FZ$. Here $U$ is an object
of $\mathcal R_{\mathrm{num}}$, whereas $\mathcal N$ is an
object of $\mathcal B$ and $F\mathcal N$ is a set.
For a recursive subset $A\subseteq\mathbb N$, let $p_A$ be
the least index of its Boolean characteristic function, with
output $1$ on $A$. Define $\theta_p(v)=M_1(\overline p,v)$.
The map $a(p)=\ulcorner\theta_p\urcorner$ from program indices
to formula codes is primitive recursive, and
$g_A:U\to U$, given by $g_A(u)=a(p_A)$, is a constant total
recursive arrow.

If $n\in A$, the equality for $Q_1$ gives
$T\vdash\theta_{p_A}(\overline n)$.
If $n\notin A$, the equality for $Q_0$ gives
$T\vdash M_0(\overline{p_A},\overline n)$, and determinism
gives $T\vdash\neg\theta_{p_A}(\overline n)$.
Consistency excludes a proof of the opposite instance in each
case. With $c_{g_A}=s(g_A\otimes1_U)$, the fixed substitution
arrow satisfies
$\sigma_{s(a(p_A),n)}=\theta_{p_A}(\overline n)$ and hence the
representation equations
\[
 c_{g_A}^*a_T^+=\pi_2^*e_A,
 \qquad
 c_{g_A}^*a_T^-=\pi_2^*e_{\mathbb N\setminus A}.
\]
The embedding of restriction idempotents for these predicates is $j_X=1_{E_X}$.

The objects $E_X$, $C_U$, and
$G=\mathsf T(U,U)=\operatorname{Tot}(\mathcal R_{\mathrm{num}})(U,U)$
belong to $\mathbf{Set}$.
Pairs of complementary restriction idempotents are precisely
$(e_A,e_{\mathbb N\setminus A})$ with $A$ recursive.
The assignment
$\rho(e_A,e_{\mathbb N\setminus A})=g_A$ is therefore a
morphism $\rho:C_U\to G$ in $\mathbf{Set}$.
The representation equations give $\ell\rho=b$ in
\eqref{eq:v8-representation-section}; the pullback
\eqref{eq:v8-representation-pullback} consequently supplies its
section $\sigma:C_U\to R$.
Both $\mathsf{Th}_T$ and $\mathsf{Rf}_T$ are c.e., so their
partial identities are arrows of $\mathcal R_{\mathrm{num}}$.
The remaining enriched category, lattice, and indexing structure
is that of Corollary~\ref{cor:numerical-definability-algebras}.
\end{proof}

\begin{corollary}\label{cor:introspection-realized-incompleteness}
For a consistent c.e. extension $T=\mathrm{PA}+\Gamma$, the
theorem and refutation sets $\mathsf{Th}_T$ and $\mathsf{Rf}_T$
are recursively inseparable.
Every consistent c.e. extension of $T$ is incomplete.
No Boolean quotient of the associated Lindenbaum algebra induces
a consistent complete presentation whose fibers over the bounds
arise from restriction idempotents.
\end{corollary}
\begin{proof}
The numerical recursion category has $1_U\ne0_{U,U}$.
Theorem~\ref{thm:introspection-numerical-definability} and
Corollary~\ref{cor:definability-algebra-inseparability} therefore
exclude a complemented restriction idempotent separating
$a_T^+$ and $a_T^-$. Such an idempotent is the partial identity
on a recursive set containing $\mathsf{Th}_T$ and disjoint
from $\mathsf{Rf}_T$.

If a consistent c.e. extension $T'$ of $T$ were complete,
simultaneous enumeration of its proofs of $\sigma_b$ and
$\neg\sigma_b$ would decide $\mathsf{Th}_{T'}$.
This recursive set contains $\mathsf{Th}_T$ and is disjoint
from $\mathsf{Rf}_T$, contradicting inseparability.

The embedding $j$ of
Lemma~\ref{lem:integration-numerical-subobjects} and
Proposition~\ref{prop:algebra-comparison} extend the same
representation section to $P_X=\operatorname{Sub}_{\mathcal P}(JX)$,
with distinguished predicates $p_T^\pm=j_Ua_T^\pm$.
Their Boolean presentation is $q_T:JU\twoheadrightarrow B_T$.
Theorem~\ref{thm:categorical-essential-incompleteness} gives the
last assertion; Proposition~\ref{prop:numerical-boolean-extensions}
identifies it with essential incompleteness of $T$.
\end{proof}

Theorem~\ref{thm:introspection-numerical-definability} derives the
representation section from arithmetic introspection after
realization in $\mathbf{Set}$. On that realized base,
Corollary~\ref{cor:numerical-introspection-separation} makes every
introspective interpretation naturally isomorphic to a constant
functor with terminal value. Thus the realized definability algebra
proves inseparability on a base that admits no nonconstant
introspective interpretation.

\appendix
\section{Definability and introspection on a fixed base}
\label{sec:introspection-representation-appendix}

Consider a finitely complete category $\mathscr V$ and a
$\mathscr V$-enriched restriction category $\mathcal K$ with
chosen binary restriction products, zeros, and bounded distributive
lattices of restriction idempotents whose joins are stable under inverse image.
Assume that its underlying ordinary category $\mathcal K_0$ is a nontrivial
recursion category, with distinguished Turing object $U$ and Turing
morphism $\tau:U\otimes U\to U$.
The Turing axiom requires a total index for each ordinary arrow.
An internal section of $\operatorname{ev}_\tau$ is a further
condition, considered below.

For the numerical enrichment over $\mathbf{Set}$,
Corollary~\ref{cor:numerical-introspection-separation} supplies a
definability algebra while every introspective functor on the base
is naturally isomorphic to a constant functor with terminal value.
The PA syntactic category admits nonconstant introspection but
no internal representation operation for the enrichment
constructed below.

\subsection{The enriched PA program category}

Let $\mathcal B$ be the exact $\Sigma_1$ syntactic category of PA
defined in Section~\ref{sec:quotation-realization}.
Write $\mathcal N$ for its equality PER defined by $x=y$,
$\bar e$ for the arithmetic numeral for $e$, and
$\underline e:1\to\mathcal N$ for the point represented by
that numeral. We use the register-program
presentation of \cite[Section~7]{RecursionCompletions}.
Write $\varphi_e$ for the partial recursive function of program
index $e$, and $\Gamma_e^{m,n}(\vec x,\vec y)$ for its graph formula
with $m$ inputs and $n$ outputs, using the fixed recursive tuple
codings. Let
\[
 \operatorname{Eq}_{m,n}(e,f)
 :=\forall\vec x\,\forall\vec y\,
   \bigl(\Gamma_e^{m,n}(\vec x,\vec y)
             \leftrightarrow\Gamma_f^{m,n}(\vec x,\vec y)\bigr).
\]
Define $H(m,n)$ to be the object of $\mathcal B$ represented by
the $\Sigma_1$ PER
$Q_{m,n}(e,f)=\operatorname{Pr}_{\mathrm{PA}}
(\ulcorner\operatorname{Eq}_{m,n}(\dot e,\dot f)\urcorner)$.
Here the dotted parameters denote numeral substitution into the
displayed formula. Primitive recursive transformations of proofs
show in PA that $Q_{m,n}$ is reflexive, symmetric, and transitive.

\begin{proposition}\label{prop:pa-enriched-recursion-category}
There is a $\mathcal B$-enriched restriction category $\mathcal K$
with objects $U_m$, $m\ge1$, hom-objects $H(m,n)$, and chosen
products $U_m\otimes U_n=U_{m+n}$. Its objects of restriction idempotents are bounded
distributive lattice objects with inverse-image-stable joins. Its underlying
ordinary category $\mathcal K_0$ is the nontrivial recursion category
$S_{\mathrm{PA}}$ of programs modulo uniformly PA-provable graph
equality. The category $\mathcal B$ admits a nonconstant
introspective structure.
\end{proposition}
\begin{proof}
The program constructors for composition and restriction induce
arrows $H(n,p)\times H(m,n)\to H(m,p)$ and
$H(m,n)\to H(m,m)$. The near-product constructor, together
with the tuple codings, induces
$H(m,n)\times H(p,q)\to H(m+p,n+q)$.
Identity, projection, diagonal, and nowhere-defined programs
give the remaining global structure arrows.
Their graph identities are provable uniformly in Heyting arithmetic
and hence in PA \cite[Lemmas~7.1--7.3]{RecursionCompletions}.

These constructors respect the PERs $Q_{m,n}$. For example,
proofs of $\operatorname{Eq}_{m,n}(e,e')$ and
$\operatorname{Eq}_{n,p}(d,d')$ can be combined with the composition
graph identity to give a proof of
$\operatorname{Eq}_{m,p}(\operatorname{cmp}(d,e),
\operatorname{cmp}(d',e'))$. The transformation of proof codes
is primitive recursive, and PA verifies it. The same construction
for the other operations proves that they descend to arrows of
$\mathcal B$. Instantiating the uniformly proved graph identities
gives proof codes for the category, restriction, and product laws;
hence these laws hold as equations in $\mathcal B$.

The object of restriction idempotents $E_{U_m}$ is the equalizer of restriction and
identity on $H(m,m)$. The join of two restriction idempotents is represented by
the partial identity on
$\operatorname{dom}(e,\vec x)\vee
\operatorname{dom}(f,\vec x)$, where
$\operatorname{dom}(e,\vec x)$ abbreviates
$\exists\vec y\,\Gamma_e^{m,m}(\vec x,\vec y)$.
The verified search construction
\cite[Lemma~7.6]{RecursionCompletions}, followed by restriction,
gives this constructor uniformly in $e,f$.
The bounded distributive-lattice identities follow from the
corresponding identities of conjunction and disjunction.
The graph formulas for inverse image prove stability of joins.
The same proof transformations make these equations internal
to $\mathcal B$.

Every global point of $H(m,n)$ is represented by a numeral
program index. Indeed, its $\Sigma_1$ graph has a PA-provable
existential output. Soundness in $\mathbb N$ supplies a standard
witness, and $\Sigma_1$-completeness proves its instance in PA.
Functionality and saturation then identify the global point
with the corresponding constant arrow. Two indices $e,f$ give
the same global point precisely when PA proves $Q_{m,n}(e,f)$.
By $\Sigma_1$-soundness and verification of finite proofs, this
is equivalent to PA proving $\operatorname{Eq}_{m,n}(e,f)$.
This identifies $\mathcal K_0$ with $S_{\mathrm{PA}}$.
Write $[e]_{m,n}:U_m\to U_n$ for the arrow represented by
program $e$, suppressing the arities when they are fixed.

The tuple codings make it an isotype. Zero programs absorb
composition and near product, so it is prodominical. The universal
program gives $\tau:U_1\otimes U_1\to U_1$.
For a representative index $e$ of an arrow
$U_m\otimes U_1\to U_1$, the verified parameter operation
produces a total primitive recursive index map on $U_m$;
its graph equation is provable in PA. This is the ordinary
Turing axiom. Consistency of PA separates identity from zero
\cite[Proposition~7.5]{RecursionCompletions}.

Take $T=\mathrm{PA}$ in the introspective structure
$(\mathcal B,\mathcal C_T,\mathsf S_T,N^T)$ of
Section~\ref{sec:quotation-realization}.
Its standard realization $F:\mathcal B\to\mathbf{Set}$ sends
$\Box_T0_{\mathcal B}$ to the empty set. If $\mathsf S_T0_{\mathcal B}$ were terminal,
$\Box_T0_{\mathcal B}=\mathcal C_T(1_{\mathcal C_T},\mathsf S_T0_{\mathcal B})$ would be terminal in
$\mathcal B$, whose image under $F$ is a singleton.
Consequently $\mathsf S_T$ is nonconstant.
\end{proof}

Here prodominicality uses zero absorption. Heller's additional
conditions for dominicality include the absence of zero divisors
\cite[Proposition~5.1]{Heller}; his definition of recursion
category uses prodominicality \cite[p.~1265]{Heller}.

\subsection{The obstruction to an internal representation operation}

Let $\operatorname{Halt}(e)$ assert that program $e$ halts on
input zero. We first record the undecidability needed for the
obstruction.

\begin{lemma}\label{lem:pa-nonhalting-undecidable}
The set $\{e:\mathrm{PA}\vdash\neg\operatorname{Halt}(\bar e)\}$
is not recursive.
\end{lemma}
\begin{proof}
Suppose a total Boolean program $d$ decides this set, with value
$1$ for membership. Write $\operatorname{par}(e,a)$ for the
verified parameter operation. Compile a binary program $q(x,z)$
which computes $d(\operatorname{par}(x,x))$, halts when the
result is $1$, and otherwise loops. If $a$ is its program index
and $b=\operatorname{par}(a,a)$, the parameter and composition
graph identities give a PA proof that program $b$ halts exactly
when $d(\bar b)$ returns $1$.

If $d(b)=1$, its finite computation record gives a PA proof
that $b$ halts. Correctness of $d$ also gives a PA proof that
$b$ does not halt, contradicting consistency.
If $d(b)=0$, its finite computation record and the loop branch
give a PA proof that $b$ does not halt, contradicting the
answer $0$. These are the two possible outputs of $d$.
\end{proof}

Fix the enrichment of
Proposition~\ref{prop:pa-enriched-recursion-category}, and put
$U=U_1$ and $G=\mathsf T(U,U)$.
The objects $E_U$, $G$, and $C_U$ are the equalizer and pullbacks
of Section~\ref{sec:generalized-definability-algebras}.

\begin{theorem}\label{thm:pa-no-internal-representation}
Equip $\mathcal K$ with predicate objects, inverse-image actions,
and monic embeddings of restriction idempotents $j_X:E_X\to P_X$ as in
Definition~\ref{def:generalized-definability-algebra}.
For every total arrow $s:U\otimes U\to S$ and predicate
$p:1\to P_S$, there is no arrow $\rho:C_U\to G$ in
$\mathcal B$ satisfying the positive representation equation
\[
 a_{U\otimes U,S}
   \langle\chi\rho,p!_{C_U}\rangle
      =\pi_2^*j_Ud_+,
 \qquad \chi(g)=s(g\otimes1_U).
\]
\end{theorem}
\begin{proof}
Let $B(e,n)$ be the primitive recursive predicate asserting that
program $e$ halts on input zero within $n$ steps.
Choose primitive recursive program constructors $a(e),a^c(e)$
whose graph identities
$\Gamma_{a(e)}(x,y)\leftrightarrow x=y\wedge B(e,x)$ and
$\Gamma_{a^c(e)}(x,y)\leftrightarrow x=y\wedge\neg B(e,x)$
are proved uniformly in PA. Their restriction and complementarity
equations have proof codes obtained uniformly by the verified
program constructors. They therefore define an arrow
$\gamma:\mathcal N\to C_U$ in $\mathcal B$.
Write $\epsilon_e:1\to E_U$ for the point representing the
program class $[a(e)]$, and
$\gamma_e=\gamma\underline e:1\to C_U$ for its complementary pair.

Suppose $\rho$ exists. A $\Sigma_1$ graph for $\rho\gamma$
allows effective search for a program index $k(e)$ representing
its output on $e$. Soundness of PA makes the search terminate,
and membership in $G$ makes the resulting function total.
Thus $c(e)=\varphi_{k(e)}(0)$ is a total recursive function.
Put $g_e=[k(e)]:U\to U$.
If $\epsilon_e=\epsilon_f$, uniqueness of complements gives
$\gamma_e=\gamma_f$, whence $g_e=g_f$ and $c(e)=c(f)$.

Conversely, suppose $c(e)=c(f)=k$. Let
$\kappa_i:U\to U$ be the total constant program with value $i$.
The two finite computations at zero yield PA proofs
$g_e\kappa_0=\kappa_k=g_f\kappa_0$.
Pull the representation equations back along
$h=\langle\kappa_0,1_U\rangle:U\to U\otimes U$.
Since $\pi_2h=1_U$, they give
\[
 j_U\epsilon_e
 =\bigl(s\langle g_e\kappa_0,1_U\rangle\bigr)^*p
 =\bigl(s\langle g_f\kappa_0,1_U\rangle\bigr)^*p
 =j_U\epsilon_f.
\]
Monicity of $j_U$ gives $\epsilon_e=\epsilon_f$.
Consequently $c(e)=c(f)$ precisely when PA proves
$\operatorname{Eq}_{1,1}(\overline{a(e)},\overline{a(f)})$.

Choose $e_0$ to be an explicit infinite loop. PA proves
$\forall n\,\neg B(\bar e_0,n)$.
The graph identities identify $\epsilon_e=\epsilon_{e_0}$ with
PA-provability of $\forall n\,\neg B(\bar e,n)$, equivalently
of $\neg\operatorname{Halt}(\bar e)$.
Comparing $c(e)$ and $c(e_0)$ therefore decides the set in
Lemma~\ref{lem:pa-nonhalting-undecidable}, a contradiction.
\end{proof}

\subsection{Internal indexing and the comparison}

The obstruction also applies to the choice of Turing indices
as an arrow of the ambient category.

\begin{proposition}\label{prop:pa-no-internal-indexing}
For the enrichment of
Proposition~\ref{prop:pa-enriched-recursion-category}, the arrow
$\operatorname{ev}_\tau:G\to H(2,1)$ has no section in
$\mathcal B$.
\end{proposition}
\begin{proof}
Suppose $\Lambda:H(2,1)\to G$ is a section.
For a unary program index $e$, write $u_e=[e]:U\to U$ for
its arrow in $S_{\mathrm{PA}}$, and put
$f_e=u_e\pi_2:U\otimes U\to U$.
The primitive recursive constructor that sends $e$ to a program
index for $f_e$ respects PA-provable equality and induces
$H(1,1)\to H(2,1)$.
Effective search in a graph for its composite with $\Lambda$
gives an index $k(e)$ of a total arrow $g_e=[k(e)]:U\to U$ satisfying
$\tau(g_e\otimes1_U)=f_e$.
As in the preceding proof, $b(e)=\varphi_{k(e)}(0)$ is total recursive.
Equality of the unary arrows represented by $e,f$ gives
$g_e=g_f$, hence $b(e)=b(f)$.

If $b(e)=b(f)=k$, finite computation gives
$g_e\kappa_0=\kappa_k=g_f\kappa_0$.
Precomposing the indexing equations with
$h=\langle\kappa_0,1_U\rangle$ therefore gives, in
$S_{\mathrm{PA}}$,
$u_e=\tau\langle\kappa_k,1_U\rangle=u_f$.
Thus $b$ decides uniformly PA-provable graph equality of unary
programs. Given $e$, compile a unary program which ignores its
input, simulates $e$ on zero, and returns zero if that simulation
halts. Its equality with the nowhere-defined program is provable
in PA exactly when $\neg\operatorname{Halt}(\bar e)$ is provable.
The resulting decision contradicts
Lemma~\ref{lem:pa-nonhalting-undecidable}.
\end{proof}

\begin{corollary}\label{cor:fixed-base-incomparability}
Among pairs consisting of a finitely complete base and an enriched
category whose underlying ordinary category is a recursion category,
as specified at the start of this appendix,
neither of the following properties implies the other:
\begin{enumerate}[label=\textup{(\roman*)}]
\item the base admits an introspective structure whose functor
is not naturally isomorphic to a constant functor with terminal value;
\item the enrichment extends to a definability algebra at its
distinguished Turing object.
\end{enumerate}
\end{corollary}
\begin{proof}
The numerical definability algebra over $\mathbf{Set}$ has
property~\textup{(ii)}, and
Corollary~\ref{cor:numerical-introspection-separation} excludes
property~\textup{(i)}.
Proposition~\ref{prop:pa-enriched-recursion-category} gives
property~\textup{(i)} for $(\mathcal B,\mathcal K)$, whereas
Theorem~\ref{thm:pa-no-internal-representation} excludes
property~\textup{(ii)}.
\end{proof}

\begin{thebibliography}{99}
\bibitem{Lawvere} F.~W.~Lawvere,
Diagonal arguments and cartesian closed categories,
in \emph{Category Theory, Homology Theory and their Applications II},
Lecture Notes in Mathematics \textbf{92}, Springer, Berlin, 1969,
134--145; reprinted with author commentary in
\emph{Reprints in Theory and Applications of Categories}
\textbf{15} (2006), 1--13.
\url{https://tac.mta.ca/tac/reprints/articles/15/tr15.pdf}.
\bibitem{Kelly} G.~M.~Kelly,
\emph{Basic Concepts of Enriched Category Theory},
London Mathematical Society Lecture Note Series \textbf{64},
Cambridge University Press, 1982; reprinted in
\emph{Reprints in Theory and Applications of Categories}
\textbf{10} (2005), 1--136.
\url{https://tac.mta.ca/tac/reprints/articles/10/tr10.pdf}.
\bibitem{CarboniCeliaMagno1982} A.~Carboni and R.~Celia Magno,
The free exact category on a left exact one,
\emph{Journal of the Australian Mathematical Society, Series A}
\textbf{33}(3) (1982), 295--301.
\url{https://doi.org/10.1017/S1446788700018735}.
\bibitem{DPH} R.~A.~Di Paola and A.~Heller,
Dominical categories: recursion theory without elements,
\emph{Journal of Symbolic Logic} \textbf{52}(3) (1987), 594--635.
\url{https://doi.org/10.2307/2274352}.
\bibitem{Heller} A.~Heller,
An existence theorem for recursion categories,
\emph{Journal of Symbolic Logic} \textbf{55}(3) (1990), 1252--1268.
\url{https://doi.org/10.2307/2274486}.
\bibitem{AspertiLongo} A.~Asperti and G.~Longo,
\emph{Categories, Types, and Structures:
An Introduction to Category Theory for the Working Computer Scientist},
Foundations of Computing, MIT Press, Cambridge, Massachusetts, 1991.
\url{https://www.cs.unibo.it/~asperti/PAPERS/book.pdf}.
\bibitem{Mendelson} E.~Mendelson,
\emph{Introduction to Mathematical Logic}, fourth edition,
Chapman \& Hall, London, 1997.
\bibitem{CarboniVitale} A.~Carboni and E.~M.~Vitale,
Regular and exact completions,
\emph{Journal of Pure and Applied Algebra} \textbf{125}(1--3) (1998),
79--116. \url{https://doi.org/10.1016/S0022-4049(96)00115-6}.
\bibitem{Bauer} A.~Bauer,
\emph{The Realizability Approach to Computable Analysis and Topology},
Ph.D. thesis, Carnegie Mellon University, 2000,
Technical Report CMU-CS-00-164.
\url{https://www.csd.cmu.edu/sites/default/files/phd-thesis/CMU-CS-00-164.pdf}.
\bibitem{RestrictionI} J.~R.~B.~Cockett and S.~Lack,
Restriction categories I: categories of partial maps,
\emph{Theoretical Computer Science} \textbf{270}(1--2) (2002),
223--259. \url{https://doi.org/10.1016/S0304-3975(00)00382-0}.
\bibitem{LambekScott} J.~Lambek and P.~Scott,
An exactification of the monoid of primitive recursive functions,
\emph{Studia Logica} \textbf{81} (2005), 1--18.
\url{https://doi.org/10.1007/s11225-005-2765-x}.
\bibitem{BBJ} G.~S.~Boolos, J.~P.~Burgess, and R.~C.~Jeffrey,
\emph{Computability and Logic}, fifth edition,
Cambridge University Press, Cambridge, 2007.
\url{https://doi.org/10.1017/CBO9780511804076}.
\bibitem{RestrictionIII} J.~R.~B.~Cockett and S.~Lack,
Restriction categories III: colimits, partial limits and extensivity,
\emph{Mathematical Structures in Computer Science}
\textbf{17}(4) (2007), 775--817.
\url{https://doi.org/10.1017/S0960129507006056}.
\bibitem{TuringCategories} J.~R.~B.~Cockett and P.~J.~W.~Hofstra,
Introduction to Turing categories,
\emph{Annals of Pure and Applied Logic} \textbf{156}(2--3) (2008),
183--209. \url{https://doi.org/10.1016/j.apal.2008.04.005}.
\bibitem{CockettSlides} J.~R.~B.~Cockett,
\emph{Turing Categories and Computability},
Estonian Winter School in Computer Science, lecture~4, March 2010.
\url{https://cspages.ucalgary.ca/~robin/talks/estonia-winter-2010/estonia-slides-4.pdf}.
\bibitem{CockettNotes} J.~R.~B.~Cockett,
\emph{Categories and Computability},
lecture notes, University of Calgary, 12 August 2014.
\url{https://cspages.ucalgary.ca/~robin/talks/estonia-winter-2010/estonia-notes.pdf}.
\bibitem{Ramesh} S.~Ramesh,
\emph{Introspective Theories and Geminal Categories},
Ph.D. dissertation, University of California, Berkeley, 2023.
\url{https://escholarship.org/uc/item/3mn0c475}.
\bibitem{RecursionCompletions} F.~Lengyel,
\emph{Completions of recursion categories},
unpublished manuscript, 22 September 2026.
\url{https://github.com/flengyel/Categorical_Incompleteness/blob/d09bd09c52c2fd398a8abbe283d9d73c79fa36d5/01_recursion_category_completions.tex}.
\bibitem{DefinabilityAlgebras} F.~Lengyel,
\emph{Definability algebras, quantifiers, and recursion categories},
unpublished manuscript, version~3, 25 October 2026.
\url{https://github.com/flengyel/Categorical_Incompleteness/blob/7abe598b9291c4bed2dddb8ae7c61de76c692428/definability_algebras_v3.tex}.
\end{thebibliography}
\end{document}
